\documentclass[11pt,a4paper]{article}
\usepackage[margin=24mm]{geometry}
\usepackage{amsmath,amssymb,amsthm,mathtools,booktabs,tabularx,microtype,xcolor,hyperref,enumitem}
\hypersetup{colorlinks=true,linkcolor=blue!45!black,urlcolor=blue!45!black,citecolor=blue!45!black}
\newtheorem{theorem}{Theorem}[section]
\newtheorem{lemma}[theorem]{Lemma}
\newtheorem{corollary}[theorem]{Corollary}
\newtheorem{proposition}[theorem]{Proposition}
\newcommand{\E}{\mathbb E}
\newcommand{\Prb}{\mathbb P}
\newcommand{\one}{\mathbf 1}
\newcommand{\TV}{\operatorname{TV}}
\newcommand{\KL}{\operatorname{KL}}
\newcommand{\Law}{\operatorname{Law}}
\newcommand{\calE}{\mathcal E}
\newcommand{\eps}{\varepsilon}
\newcommand{\sig}{\operatorname{sigmoid}}
\setlength{\emergencystretch}{2em}
\title{Direct Acquired Prevention\\Without Absolute Forecasts}
\author{Astra research continuation\\Proofs, finite protocols, execution and preserved sources}
\date{5 November 2026}
\begin{document}
\maketitle
\begin{abstract}
We construct a decision that selects a nearly optimal kinetic intervention for a rare first event without acquiring its absolute rarity, lifetime, complete generator, CDF or marked record. The principal family consists of heterogeneous interacting binary components with a common unknown Gibbs invariant law, uniformly reliable local conditionals, bounded operating clocks, and a certified mixing margin. Failure is the single all-failed configuration; this includes arbitrarily large parallel redundant systems with interacting components. The preferred experiment measures only an average of local escape rates in the target's neighborhood. An explicit Green-function and superharmonic calculation bounds its correction to the true mean and cancels the unknown rare probability from the comparison. Binary finite holds suffice. A second, grouped first-exit assay has the ordinary inverse-square accuracy cost, with finite censoring, clipping, clock and repeated-cell errors and a matching exposure lower bound. Neither branch simulates a global trajectory.

The boundary certificate also controls the mean from every safe starting configuration. A graphical comparison then makes a bounded operating warm-up sufficient, with an explicit initially-failed mass bound covering the entire finite clock window. No inspection or equilibrium preparation is required for this mean-only branch. A finite one-bit alternative and the stronger relative-equilibrium preparation are retained separately. When the boundary correction is too large, a target-start return probe stops after only logarithmically many components have recovered. Fixed-duration local experiments then implement a different reversible generator whose killed Dirichlet form uniformly brackets the operating generator. Timing accuracy for that comparison is set by the decision tolerance rather than the number of simulated updates; the other service errors retain their distinct costs. A finite full-return fallback retains polynomial cost without inverse global rarity.

At fixed degree, contrast, accuracy and confidence, coarse acquisition costs $\widetilde O(N)$ local services; the subsequent decision uses $O(K)$ boundary experiments in the certified boundary regime, or $O(K\operatorname{polylog}N)$ updates in the return regime, with accuracy and confidence factors explicit. In the same local experimental access model, absolute relative lifetime prediction requires $\Omega(N^2)$ total local information even for independent components with fast known clocks. This is a separation from forecasting, not a universal lower bound against specialist direct-decision methods. Actual pin, field, detector, timing, repeated-cell, preparation and deployment services remain priced, including the difference between whole-system wall time and summed component exposure. No economical general device or historical priority is claimed. Independent continuous completion design, exact assignment decisions, and controlled quantum export without Hamiltonian learning are retained separately.
\end{abstract}
\tableofcontents
\newpage
\section{The question and the exact endpoint}
Can an experimentally supported prevention decision cost less than the absolute prediction that one might ordinarily construct first? The answer here is affirmative in a declared family. The action is chosen before deployment from a finite feasible menu. The benchmark is the mean time to a first catastrophic state from the equilibrium-conditioned safe law. The strongest initialization theorem transfers the decision to every safe starting state and to an actual finite operating warm-up with bounded initial-failure mass, without preparing that full law. It is a relative mean objective, not an arbitrary deadline risk, terminal-damage loss, adaptive policy or certificate of model completeness.

There are $N$ observed bits and a known interaction graph of maximum degree $d$. A one means failed. Set $b=(1,\ldots,1)$, $B=\{b\}$ and $A=B^c$. In menu item $u\in\mathcal U$, $|\mathcal U|=K$,
\begin{equation}
 Q^u=\sum_{i=1}^NQ_i^u,\quad Q_i^u(x,x^i)=r_i^u(x_{\partial i})
 \begin{cases}q_i(x_{\partial i}),&x_i=0,\\1-q_i(x_{\partial i}),&x_i=1.\end{cases}
 \label{eq:model}
\end{equation}
The common unknown equilibrium law is
\begin{equation}
 \pi(x)\propto\exp\left(\sum_i h_i x_i+\sum_{ij}J_{ij}x_ix_j\right),\quad
 q_i(s)=\sig\left(h_i+\sum_jJ_{ij}s_j\right).
 \label{eq:gibbs}
\end{equation}
The menu changes reciprocal kinetic clocks, not these equilibrium weights. The clocks can be different unknown functions of context for every menu item; no known multiplier response is required. The following numerical envelopes must hold simultaneously:
\begin{equation}
 0<r_-\le r_i^u(s)\le r_+,\qquad q_i(s)\le\bar q<\tfrac12,\qquad
 \operatorname{gap}(-Q^u)\ge\gamma>0.
 \label{eq:envelopes}
\end{equation}
A certified Dobrushin margin $\alpha$ supplies $\gamma=\alpha r_-$ through the accepted heat-bath backbone comparison. Section~\ref{sec:acquisition} constructs coarse finite local evidence for the required margins. Exact common-invariant closure, faithful local copies, reciprocal dynamics, and absence of omitted mechanisms remain premises. Finite data cannot prove those premises against arbitrary alternatives.

The target is
\begin{equation}
 M_u=\E_{\pi(\cdot\mid A),Q^u}\tau_b,\qquad
 M_{\widehat u}\ge(1-\eps)\max_{u\in\mathcal U}M_u.
 \label{eq:goal}
\end{equation}
Known costs may first restrict the feasible menu. A monetary cost divided into an unknown absolute lifetime is a different objective. Preparation-density and deployment-generator errors are included later.

The configuration space has $2^N$ states, and the event can be exponentially rare. No exchangeability, birth--death reduction, independence of components or explicit global matrix is assumed. Nonuniform fields, signed weak interactions and context-dependent clocks on bounded-degree graphs are allowed. The singleton event is substantive: arbitrary all-terminal disconnection is covered by the preceding general machinery, not automatically by the direct score below.
\section{Finite local holds preserve a decision without preserving paths}
\label{sec:hold}
One local update chooses a coordinate uniformly, clamps its neighbors to the virtual context, pins its current bit, reproduces the actual operating environment, releases for a prescribed hold, and reads its final bit. It records neither an exit time nor intermediate transitions. It estimates no rate and manufactures no fine-accuracy uniformization coin.

\subsection{A general reversible-block result}
Let $Q=\sum_iQ_i$, with each finite generator $Q_i$ reversible for the same $\pi$, and spectrum of $-Q_i$ in $[0,\Lambda]$. For an exact duration $h>0$ define
\begin{equation}
 P_h=\frac1N\sum_i e^{hQ_i},\quad Q_h=\frac Nh(P_h-I),\quad
 \phi(z)=\frac{1-e^{-z}}z.
 \label{eq:holdkernel}
\end{equation}
The kernel $P_h$ is reversible and positive semidefinite in $L^2(\pi)$. Spectral calculus gives
\begin{equation}
 \phi(h\Lambda)\calE_Q(f,f)\le\calE_{Q_h}(f,f)\le\calE_Q(f,f).
 \label{eq:formhold}
\end{equation}
This uses $\phi(h\Lambda)x\le(1-e^{-hx})/h\le x$ for $0\le x\le\Lambda$, not a trajectory approximation.

\begin{theorem}[Finite holds and first-event means]\label{thm:hold}
For any nonempty proper target $B$ in an irreducible finite reversible chain, put $A=B^c$. If a second generator sharing $\pi$ obeys
$\kappa_-\calE_Q\le\calE_{\widetilde Q}\le\kappa_+\calE_Q$, where $0<\kappa_-\le\kappa_+$, then
\begin{equation}
 \frac{M_Q}{\kappa_+}\le M_{\widetilde Q}\le\frac{M_Q}{\kappa_-},\qquad
 M_Q=\E_{\pi(\cdot\mid A),Q}\tau_B.
 \label{eq:meanhold}
\end{equation}
There is no rarity, gap, state-count or path-length factor in this transfer.
\end{theorem}
\begin{proof}
On functions zero on $B$, $K=-Q_A$ and $\widetilde K=-\widetilde Q_A$ are positive definite in $L^2(\pi|_A)$. Their form inequality reverses under inversion. The Poisson equation yields
$M_Q=\langle\one_A,K^{-1}\one_A\rangle_\pi/\pi(A)$.
Apply the inverse inequalities. The reversible variational structure is classical; its use here certifies an experimentally implementable substitute chain.
\end{proof}

\subsection{Actual finite timing errors}
In a binary context $s$, $Q_i^u=r_i^u(s)(H_i-I)$, where $H_i$ redraws the bit with probability $q_i(s)$. A hold of duration $t_i^u(s)$ gives
\begin{equation}
 e^{t_i^u(s)Q_i^u}=I+(1-e^{-r_i^u(s)t_i^u(s)})(H_i-I).
 \label{eq:localhold}
\end{equation}
The actual hold may be unknown in $h[1-\omega,1+\omega]$. At a fixed context it must be shared by the two possible focal starting bits, and define a fixed kernel throughout the decision. A fixed conditional distribution of durations, shared by the two bits, also works. Arbitrary history-dependent timing does not supply a single reversible kernel.

Set
\begin{equation}
 \kappa_-=(1-\omega)\phi(r_+h(1+\omega)),\qquad \kappa_+=1+\omega.
 \label{eq:kappa}
\end{equation}
The actual ideal generator is $Q_h=\sum_i a_i^u(s)Q_i^u$, with
$a_i^u(s)=[1-e^{-r_i^u(s)t_i^u(s)}]/[r_i^u(s)h]\in[\kappa_-,\kappa_+]$.
Thus reversibility and the common invariant law are exact for this unknown substitute generator. The timing uncertainty does not accumulate with the number of updates. For example,
\begin{equation}
 h=\frac{\eps}{8r_+},\quad \omega\le\frac{\eps}{32},\quad
 \frac{\kappa_-}{\kappa_+}\ge1-\frac\eps4\quad(0<\eps\le1/4)
 \label{eq:holdrecipe}
\end{equation}
follows conservatively from $\phi(z)\ge1-z/2$. Its gap is at least
\begin{equation}
 g=\min\{1,h\kappa_-\gamma/N\}.
 \label{eq:g}
\end{equation}
Only the selected local context is realized physically; the rest is classical memory. Intermediate physical flips are intentionally omitted. The implemented object is $P_h$, whose induced generator is compared to $Q$, not falsely equated to it.

Pin failures, incorrect final reads, context errors and uncontrolled transients need not preserve $\pi$. They are budgeted separately below. Nor does this theorem permit clocks measured at compensating fields to be transported to operation. It covers all event sets for the equilibrium-safe mean, but asserts neither marked-path TV closeness nor pointwise ordering of every first-hit CDF.

\section{A short return score cancels the unknown event rarity}
\label{sec:return}
The remainder uses the singleton $b$. Replacing it with an arbitrary event set would generally require a conditional distribution and a matrix-valued return interface; no such extension is assumed.

\subsection{The exact canceled quantity}
For $P=P_h^u$, $p=\pi(b)$, define
\begin{equation}
 Z(x)=\sum_{j\ge0}(P^j(x,b)-p),\qquad S_u=Z(b).
 \label{eq:green}
\end{equation}
The Poisson equation $(I-P)Z=\one_{\{b\}}-p$ and $\sum_x\pi(x)Z(x)=0$ imply
\begin{equation}
 \E_x T_b=\frac{S_u-Z(x)}p,\qquad
 M_{h,u}=\frac hN\frac{S_u}{p(1-p)}.
 \label{eq:cancel}
\end{equation}
The enormous unknown factor $1/[p(1-p)]$ is exactly common to all items. Ranking $S_u$ ranks their substitute means. This identity is established Green-function theory; the finite experimental reduction is the operation constructed here.

Since full conditionals are at most $\bar q$, the same is true under partial conditioning. Multiplying a chain of conditional probabilities, and separately multiplying odds along bit flips, gives
\begin{equation}
 p\le\bar p:=\bar q^N,\qquad
 \frac{\pi(b)}{\pi(x)}\le a^{w(x)},\quad
 a=\frac{\bar q}{1-\bar q}<1,\quad w(x)=\#\{i:x_i=0\}.
 \label{eq:massratio}
\end{equation}
No partition function or numerical weight table is required. Put
\begin{equation}
 d_-=1-e^{-r_-h(1-\omega)},\quad d_+=1-e^{-r_+h(1+\omega)},\quad
 s_0=\frac{(1-\bar p)^2}{d_+}.
 \label{eq:s0}
\end{equation}
The PSD spectral representation and Cauchy--Schwarz give
\begin{equation}
 S_u\ge\frac{(1-p)^2}{1-P(b,b)}\ge s_0.
 \label{eq:greenlower}
\end{equation}
The nonconstant spectral weights sum to $1-p$, their weighted $1-\lambda$ sum is $1-P(b,b)$, and the latter is at most $d_+$.

\subsection{Stop after a few recoveries}
For $1\le k<N$ let $C_k=\{x:w(x)\ge k\}$ and let $\sigma_k$ be its first entrance time from $b$. Count
\begin{equation}
 V_u=\sum_{j=0}^{\sigma_k-1}\one_{\{X_j=b\}},\qquad A_u=\E_b V_u.
 \label{eq:score}
\end{equation}
Each visit to $b$, including a self-loop, restarts the same chance of reaching $C_k$ before the next visit. Therefore $V_u$ is geometric on $\{1,2,\ldots\}$ and
\begin{equation}
 A_u=\frac1{\Pr_b(\sigma_k<T_b^+)},\qquad
 \operatorname{Var}(V_u)=A_u(A_u-1)\le A_u^2.
 \label{eq:geometric}
\end{equation}
Equivalently $A_u=\pi(b)/\operatorname{cap}_P(b,C_k)$. Capacity and renewal theory supply the representation; the next bounds make this unknown-system measurement finite and short.

While $w<k$ its drift is at least
\begin{equation}
 \nu=\left(1-\frac kN\right)d_-(1-\bar q)-\frac kN d_+\bar q.
 \label{eq:drift}
\end{equation}
Require $\nu>0$. Optional stopping and the bounded-increment martingale inequality yield
\begin{equation}
 \E_b\sigma_k\le k/\nu,\qquad
 \Pr_b(\sigma_k>J)\le e^{-(\nu J-k)^2/(2J)},\quad J\ge k/\nu.
 \label{eq:stopcost}
\end{equation}
For the second assertion extend the stopped count by deterministic drift $\nu$ after exit. Its increments have range length two and conditional drift at least $\nu$. On $\sigma_k>J$ it is still below $k$. If $J\ge2k/\nu$ the displayed bound is at most $e^{-\nu^2J/8}$.

\begin{theorem}[Finite local-return certificate]\label{thm:localreturn}
For any positive integer analysis horizon $L$, set
\begin{equation}
 D=\bar p(L+k/\nu)+La^k+\frac{a^{k/2}e^{-gL}}g.
 \label{eq:D}
\end{equation}
Then, uniformly over the menu, $|S_u-A_u|\le D$. The horizon $L$ is used only in certificate arithmetic, not as the running length of the probe.
\end{theorem}
\begin{proof}
For a boundary state $x$ with $w(x)=k$, reversibility gives $P^j(x,b)\le\pi(b)/\pi(x)\le a^k$. The first $L$ terms of $Z(x)$ are between $-Lp$ and $La^k$. Spectral contraction gives
\[
 |P^j(x,b)-p|\le\sqrt{\pi(b)/\pi(x)}(1-g)^j\le a^{k/2}(1-g)^j,
\]
so the remaining sum has modulus at most $a^{k/2}e^{-gL}/g$. The stopped Poisson identity is
$S_u=A_u-p\E_b\sigma_k+\E_bZ(X_{\sigma_k})$.
Apply these bounds and~\eqref{eq:stopcost}. The finite state space and positive drift justify stopping.
\end{proof}

For $0<\eta_b<1/4$, one finite schedule is
\begin{align}
 L&=\max\left\{1,\left\lceil g^{-1}\log\max\left(1,\frac3{\eta_bgs_0}\right)\right\rceil\right\},\\
 k&=\max\left\{1,\left\lceil\frac{\log\max(1,3L/(\eta_bs_0))}{\log(1/a)}\right\rceil\right\}.
 \label{eq:schedule}
\end{align}
Check $k<N$, $\nu>0$, $\bar p(L+k/\nu)\le\eta_b s_0/3$, and directly check $D\le\eta_b s_0$ with outward rounding. At fixed margins these checks succeed for all sufficiently large $N$: $g^{-1}=O(N)$, $L=O(N\log N)$, $k=O(\log N)$, and $\bar p$ decreases exponentially. The threshold is computed. A boundary without certified slack triggers refinement or the fallback below.

\subsection{Confidence, finite caps and ranking}
A geometric variable of mean $A\ge1$ obeys
$\E e^{tV}\le(1-At)^{-1}$ for $0\le t<1/A$ and $\E V^2\le2A^2$.
For the first bound use its exact generating function and $(1-t)e^t\le1$; for the lower tail use $e^{-x}\le1-x+x^2/2$. Chernoff's argument gives
\begin{equation}
 \Pr\big(|\widehat A_u/A_u-1|>\eta_s\big)\le2e^{-n\eta_s^2/4},
 \quad0<\eta_s\le1.
 \label{eq:concentration}
\end{equation}
The conditional version follows by iteration of conditional moment bounds; no independence of reused physical cells is presumed beyond the declared conditional interface.
Use
\begin{equation}
 n=\left\lceil4\eta_s^{-2}\log\frac{2K}{\delta_{\rm stat}}\right\rceil,
 \quad J=\left\lceil\max\left\{2k/\nu,\,8\nu^{-2}\log\frac{Kn}{\delta_{\rm cap}}\right\}\right\rceil.
 \label{eq:finitecaps}
\end{equation}
Abort explicitly if any probe fails to reach $C_k$ by $J$ updates. Its ideal probability is at most $\delta_{\rm cap}$. Otherwise the finite algorithm agrees with the uncapped reference. Choosing the largest sample mean then gives, with confidence $1-\delta_{\rm stat}-\delta_{\rm cap}$,
\begin{equation}
 \frac{M_{\widehat u}}{\max_uM_u}
 \ge\frac{\kappa_-}{\kappa_+}
 \frac{(1-\eta_s)(1-\eta_b)}{(1+\eta_s)(1+\eta_b)}.
 \label{eq:decisionfactor}
\end{equation}
Here $D\le\eta_b s_0\le\eta_b S_u$ and the common $p$ cancels.

The observed certificate can be sharper. Set
\begin{equation}
 S_u^- =\max\{s_0,\widehat A_u/(1+\eta_s)-D\},\quad
 S_u^+ =\widehat A_u/(1-\eta_s)+D.
 \label{eq:observedinterval}
\end{equation}
Then $M_u/M_v\ge(\kappa_-/\kappa_+)S_u^-/S_v^+$. If the latter exceeds one for every competitor, the protocol certifies the actual best mean-lifetime item while leaving the absolute lifetime unknown. Adaptive sample stages must divide their confidence budget; repeated use of a fixed-time interval is invalid.

\subsection{A finite fallback without rarity reconstruction}
If the sparse certificate fails, use full length-$L$ counts $V_L=\sum_{j<L}\one_{\{X_j=b\}}$. Their means are $S_u+Lp-T_u$, where $0\le T_u\le e^{-gL}/g$. Choose $L$ so that this last bound is at most $\eta s_0$, and estimate each count mean to additive $\eta s_0$ using
\[
 n_{\rm fb}=\left\lceil\frac{L^2}{2\eta^2s_0^2}\log\frac{2K}{\delta_{\rm stat}}\right\rceil.
\]
The common unknown $Lp$ cancels in comparisons. The selected $S$ is at least the best minus $3\eta s_0$, hence within relative factor $1-3\eta$. This is polynomial in $N$, numerical inverse margins and accuracy, not in $1/p$. At small $N$ an ordinary exact solver can be better; this fallback is not claimed optimal.

\section{The preferred branch: a direct boundary experiment}
\label{sec:boundary}
The short-return method is not the cheapest applicable method at fixed accuracy in a large uniformly reliable system. A boundary-flux approximation can be certified uniformly. Incorporating it removes trajectory simulation entirely. This comparison changed the implementation without changing the prevention objective.

\subsection{A proved correction to the ordinary boundary-flux approximation}
For analysis only define $R=Nr_+$ and $P^\circ=I+Q^u/R$. It is never implemented. Each $-Q_i^u\preceq r_+I$, so $P^\circ$ is positive semidefinite; its gap is at least $g_0=\gamma/(Nr_+)$. Define $S_u^\circ=\sum_{j\ge0}[(P^\circ)^j(b,b)-p]$ and
\begin{equation}
 q_{0,u}=1-P^\circ(b,b)=\frac{Q^u(b,B^c)}{Nr_+}.
 \label{eq:q0}
\end{equation}
Then $M_u=S_u^\circ/[R p(1-p)]$ exactly. Let
\begin{equation}
 C=\frac{r_+\bar q}{r_-(1-\bar q)},\quad
 \rho=\frac{Ck}{N-k}<1,\quad z=\frac{C}{(N-1)(1-\rho)}<1.
 \label{eq:boundaryconstants}
\end{equation}
After leaving $b$, the probability of returning before reaching $C_k$ is at most $z$. This does not require the Hamming weight to be Markov. At any configuration with $w$ working bits, its downward/upward move ratio is at most $Cw/(N-w)$. Put $r_j=Cj/(N-j)$, $R_0=1$, $R_j=\prod_{l=1}^j r_l$, and
\[
 H(w)=\frac{\sum_{j=w}^{k-1}R_j}{\sum_{j=0}^{k-1}R_j},\qquad H(0)=1,\ H(k)=0.
\]
Its conditional drift is nonpositive for every actual configuration, since successive increments have ratio $r_w$. Optional stopping bounds the return probability by $H(1)\le\sum_{j\ge1}r_1\rho^{j-1}=z$. This is a birth--death superharmonic comparison, not an exact lumping assertion.

\begin{theorem}[Uniform boundary certificate]\label{thm:boundary}
For any positive integer analysis horizon $L$ and $k$ satisfying~\eqref{eq:boundaryconstants}, set
\begin{equation}
 \ell=(1-\bar p)^2,\qquad
 U=\frac1{1-z}+La^k+\frac{a^{k/2}e^{-g_0L}}{g_0}.
 \label{eq:boundaryenvelope}
\end{equation}
Then, for every menu item,
\begin{equation}
 \boxed{\ell\le q_{0,u}S_u^\circ\le U.}
 \label{eq:boundarybridge}
\end{equation}
\end{theorem}
\begin{proof}
The lower bound is $(1-p)^2\le q_{0,u}S_u^\circ$. For the upper bound use the stopped Poisson identity $S_u^\circ=A_u-p\E\sigma_k+\E Z(X_{\sigma_k})$. The return count is geometric with mean $A_u\le[q_{0,u}(1-z)]^{-1}$. Drop the nonpositive middle term, apply the upper boundary bound on $Z$ from Theorem~\ref{thm:localreturn}, and multiply by $q_{0,u}\le1$.
\end{proof}
At fixed positive local margins choose $L=O(N\log(N/\eta))$ and $k=O(\log(N/\eta))$ so that the last two terms are at most $\eta$. Then $z=O(1/N)$, $\ell\to1$, and $U\le1+O(1/N)+\eta$. Every inequality is explicitly checked. No uncontrolled rare-event or Poisson-clumping approximation is used.

\subsection{A binary local endpoint measurement}
Let
\[
 \lambda_i^u=Q_i^u(b,b^i)=r_i^u(\one)(1-q_i(\one)),\quad
 \overline\lambda_u=N^{-1}\sum_i\lambda_i^u,\quad\theta_u=h\overline\lambda_u.
\]
Only the all-failed neighbor context is queried. A uniform random coordinate is pinned to one at its actual operating menu. After a nominal hold $h$, record whether it is zero. Equilibrium probabilities and clocks are never estimated separately.

If the actual hold $T\in[h(1-\omega),h(1+\omega)]$ is exogenous to the focal path conditional on prior history, its success probability obeys
\begin{equation}
 \kappa_-\theta_u\le\mu_{u,j}\le\kappa_+\theta_u,
 \quad\kappa_-=(1-\omega)\phi(r_+h(1+\omega)),\quad\kappa_+=1+\omega.
 \label{eq:directresponse}
\end{equation}
The duration may change with previous observations. No fixed experimental kernel and no symmetry of timings for two starting bits is required: this branch starts only at one.

If the timer responds to the current focal path, use the weaker
\begin{equation}
 \kappa_-^{\rm stop}=(1-\omega)e^{-r_+h(1+\omega)}.
 \label{eq:stoppingtimer}
\end{equation}
One recovery before the lower deadline followed by no transition before the upper deadline forces every admissible stopped read to be zero. Its probability is at least $\lambda_i^u h(1-\omega)e^{-r_+h(1+\omega)}$. Success also requires an exit before the upper deadline, with probability at most $\lambda_i^u h(1+\omega)$. The weaker coefficient is not silently replaced by the exogenous one.

Let the combined conditional mean error of pin, context, selector, read and arithmetic services be at most $\beta_0$. The response bounds acquire additive $\pm\beta_0$. Since $\theta_u\ge\theta_{\min}=h r_-(1-\bar q)$, setting $\beta_0\le\eta_b\kappa_-\theta_{\min}$ gives
\[
 a_-\theta_u\le\E[Y_{u,j}\mid\mathcal F_{j-1}]\le a_+\theta_u,
 \quad a_-=(1-\eta_b)\kappa_-,\ a_+=(1+\eta_b)\kappa_+.
\]
This allowance is set by the decision, not $N$, $1/p$ or the number of trials. Dependence of reused cells consistent with these conditional envelopes is allowed.

\begin{theorem}[Acquired boundary decision]\label{thm:decision}
Take
\begin{equation}
 n\ge\frac3{\eta_s^2a_-\theta_{\min}}\log\frac{2K}{\delta_{\rm dec}},
 \qquad0<\eta_s\le1,
 \label{eq:boundarysamples}
\end{equation}
local binary trials per item and choose the smallest recovery frequency. With probability at least $1-\delta_{\rm dec}$,
\begin{equation}
 \boxed{\frac{M_{\widehat u}}{\max_uM_u}
 \ge\frac\ell U\frac{\kappa_-}{\kappa_+}
 \frac{(1-\eta_s)(1-\eta_b)}{(1+\eta_s)(1+\eta_b)}.}
 \label{eq:boundarydecision}
\end{equation}
No trajectory, event clock, static counter, rare-state sampler, learned kinetic table or marked-law certificate is used.
\end{theorem}
\begin{proof}
Conditional Bernoulli exponential bounds give simultaneous inequalities
$(1-\eta_s)a_-\theta_u\le\overline Y_u\le(1+\eta_s)a_+\theta_u$ with the stated confidence; their Chernoff exponents are at least $n\eta_s^2a_-\theta_{\min}/3$. The selected $\theta$ is at most $(1+\eta_s)a_+/[(1-\eta_s)a_-]$ times that of any competitor. Use $q_{0,u}=\theta_u/(hr_+)$ and~\eqref{eq:boundarybridge}. The common $p(1-p)$ cancels exactly.
\end{proof}
For any $0<\eps\le1/4$, use~\eqref{eq:holdrecipe}, choose $\eta_s,\eta_b$ as sufficiently small fixed multiples of $\eps$, and require $\ell/U\ge1-c\eps$ with remaining explicit slack for preparation and deployment. The check succeeds for sufficiently large $N$ at fixed margins. If it fails, the return protocol resolves the same objective; an uncontrolled approximation is not substituted.

\subsection{Observable finite-bit certificates}
The implementation uses tighter one-sided Bernoulli Chernoff inversion. For $s$ successes in $n$ trials, an endpoint $p$ is certified by $e^{-n\operatorname{kl}(s/n,p)}\le\delta_{\rm tail}$. At dyadic $p=A/D$ this is an integer comparison:
\[
 A^s(D-A)^{n-s}n^n\delta_{\rm tail}^{-1}
 \le D^n s^s(n-s)^{n-s},
\]
with rational confidence denominators multiplied across. Outward binary bisection gives $[m_u^-,m_u^+]$ with the two conditional-mean guarantees. Therefore
\begin{equation}
 \theta_u^- =\max\{0,(m_u^- -\beta_0)/\kappa_+\},\quad
 \theta_u^+ =(m_u^+ +\beta_0)/\kappa_-,\qquad
 \frac{M_u}{M_v}\ge\frac\ell U\frac{\theta_v^-}{\theta_u^+}.
 \label{eq:observedboundary}
\end{equation}
This is valid for dependent trials obeying the conditional bounds; their physical sampling law need not be binomial.

For structural arithmetic the code uses rational parameters, $e^{-g_0L}\le2^{-\lfloor g_0L\rfloor}$ and $a^{k/2}\le a^{\lfloor k/2\rfloor}$. Coefficients $\phi$ or $e^{-z}$ for $0\le z\le1$ are bounded by alternating rational Taylor sums. Every issuing comparison is exact rational or integer arithmetic. Displayed decimals are rounded values.

\providecommand{\RAE}{\mathbb E}
\providecommand{\RAP}{\mathbb P}
\providecommand{\RAone}{\mathbf1}
\providecommand{\RAKL}{\operatorname{KL}}

\section{A finite direct assay for the mean boundary escape rate}
\label{sec:rate_assay}
The required decision statistic is an arithmetic average of actual local
operating rates. For a menu setting $u$ and target configuration $b$, let
\[
 \lambda_i^u=Q^u(b,b^{(i)}),\qquad
 \mu_u=\frac1N\sum_{i=1}^N\lambda_i^u,
 \qquad \Lambda_u=N\mu_u.
\]
Suppose a common certified envelope supplies
\begin{equation}
 0<\lambda_-\le\lambda_i^u\le\lambda_+<\infty,
 \qquad \kappa=\lambda_+/\lambda_-.
 \label{ra:bounds}
\end{equation}
For the preserved binary Gibbs kinetics with $b_i=1$,
$\lambda_i^u=r_i^u(b_{\partial i})[1-q_i^u(b_{\partial i})]$.
Thus $r_i^u\ge r_-$ and $q_i^u\le\bar q<1$ give the usable fast-direction
floor $\lambda_-=r_-(1-\bar q)$, and $r_i^u\le r_+$ gives
$\lambda_+=r_+$. These are bounds on the requested operating environments,
not rates transported from balancing fields.
Random-context averaging does not establish a uniform bound on unvisited
contexts. A supplied engineering bound or an independently acquired
uniform certificate retains its own full acquisition cost.

For each trial choose a coordinate uniformly, install its declared
neighborhood in the target context on a faithful local copy, and make
four successive pin--release first-exit observations from the same
requested starting bit. In the ideal Markov model the four waits at
coordinate $i$ are independent $\operatorname{Exp}(\lambda_i^u)$
conditional on the requested coordinate. Repeated physical cells need
only the uniform conditional implementation guarantees below. No
independence between imperfect trials or cells is assumed.

The context setting, four starting-state preparations, settling,
operating-field reproduction, detector operation and clock service are
all paid. In a bounded-neighborhood model this sets only the declared
local neighborhood; it does not prepare the entire global target state.
A local pin is not a free natural visit to that state. Structural closure
and faithful local replication remain premises.

\subsection{The classical inverse-gamma statistic and finite modifications}
Let $T_1,\ldots,T_4$ have rate $\lambda$ and set $S=\sum_jT_j$.
The classical unbiased rate statistic is $X=3/S$. Direct integration of
the gamma density gives
\begin{equation}
 f_S(s)=\frac{\lambda^4s^3e^{-\lambda s}}6,
 \qquad \RAE X=\lambda,
 \qquad \RAE X^2=\frac32\lambda^2.
 \label{ra:gamma}
\end{equation}
This identity is standard exponential-family estimation, not an invented
estimator; it is derived here to expose the exact variance and finite
device requirements. In particular, pooling heterogeneous waits and
inverting their overall mean would estimate a harmonic-rate quantity,
not the required arithmetic mean. The coordinate is held fixed within
each group of four, then the resulting rate statistics are averaged.

Fix $0<\eta\le1/4$. Choose a rational $a$ with
\begin{equation}
 \frac{\eta^{1/3}}{2\lambda_+}<a\le
 \frac{\eta^{1/3}}{\lambda_+},\qquad B=3/a,
 \qquad D=\frac{\lceil\log_2(192/\eta)\rceil}{\lambda_-}.
 \label{ra:clipcap}
\end{equation}
A dyadic $a\lambda_+$ satisfying these inequalities is obtained by
integer comparisons of cubes. Let
\[
 X_a=\min(3/S,B)=3/\max(S,a),\qquad
 Z=X_a\RAone_{\{\max_jT_j\le D\}}.
\]
Run the waits sequentially, stopping and returning zero as soon as one
wait is censored. The ideal output lies in $[0,B]$. The two lost-mean
bounds are
\begin{align}
 0\le\lambda-\RAE X_a
 &\le\frac{\lambda^4a^3}{24}
 \le\frac{\eta\lambda}{24},\label{ra:clipbias}\\
 0\le\RAE X_a-\RAE Z
 &\le\frac{12e^{-\lambda D}}D
 \le\frac{\eta\lambda_-}{16}.
 \label{ra:censorbias}
\end{align}
For \eqref{ra:clipbias}, integrate
$(3/s-3/a)f_S(s)$ on $[0,a]$ and bound $e^{-\lambda s}$ by one.
For \eqref{ra:censorbias}, on any $T_j>D$ one has $X_a\le3/D$;
take a union bound over four waits. Also $D\lambda_-\ge1$ and
$e^{-\lambda_-D}\le2^{-\lambda_-D}\le\eta/192$.
The clipped, censored second moment is at most $3\lambda^2/2$.

\subsection{Actual deadlines, first-event timestamps and finite arithmetic}
An actual timer must supply the following service, including start,
release, detector latency and deadline response in the error budget:
outside a declared bad-service event, an uncensored first-exit timestamp
differs from the true wait by at most $\Delta$, and each actual censor
deadline differs from $D$ by at most $\Delta$. Take
\begin{equation}
 0<\Delta\le\frac{\eta\lambda_-a^2}{1024},\qquad
 0<\omega\le\frac{\eta\lambda_-}{128},\qquad
 0\le\beta\le\frac{\eta\lambda_-}{128B}.
 \label{ra:hardware}
\end{equation}
$\omega$ bounds final numerical rounding of the rate statistic. The
total conditional TV or bad-service probability allowance $\beta$ per
trial covers finite coordinate selection, context and pin failures,
missed exits, nonideal local conditional laws and other declared
stochastic readout errors. It is a guarantee conditional on the entire
past, on the common calibration good event. Deterministic timer errors
are accounted for separately and are not represented as a false TV
approximation of a continuous timestamp by a grid.

The reported statistic is
$Y=3/\max(a,\sum_j\widehat T_j)$ if all four observations are accepted,
and zero on a censor or flagged failure, rounded within $\omega$ and
clipped to $[0,B]$. Round each timestamp with its physical errors included
in $\Delta$. Exact rational division is also allowed. Every operation
has a finite cap, including failed reads and state preparations.

The censor decision can disagree with its ideal value only when at
least one true wait lies in $[D-\Delta,D+\Delta]$, or on the bad-service
event. The ideal censor-boundary band has probability at most
$8\lambda_+\Delta$. Away from that band, the reciprocal clipping map is
$3/a^2$-Lipschitz, so its time-induced error is at most $12\Delta/a^2$.
Consequently
\begin{equation}
 |\RAE(Y\mid\mathcal F)-\RAE Z|
 \le\frac{12\Delta}{a^2}+8B\lambda_+\Delta+B\beta+\omega.
 \label{ra:timerbias}
\end{equation}
Here $\mathcal F$ contains every previous trial and acquisition outcome;
the ideal expectation includes a fresh uniform coordinate.
The band term is essential: a Lipschitz bound on a recorded sum alone
would not account for censor decisions changed by timer error.

\begin{lemma}[Uniform conditional moments of the physical trial]
On the simultaneous calibration good event, the construction satisfies
\begin{equation}
 |\RAE(Y\mid\mathcal F)-\mu_u|\le\frac{\eta\mu_u}{6},
 \qquad \RAE(Y^2\mid\mathcal F)\le4\lambda_+\mu_u
                                      \le4\lambda_+^2.
 \label{ra:trialmoments}
\end{equation}
\end{lemma}
\begin{proof}
Average \eqref{ra:clipbias} over the uniform coordinate. The clipping
loss is at most $\eta\mu_u/24$. The censor loss is at most
$\eta\lambda_-/16$. Since $\lambda_+a\le1$, the remaining errors in
\eqref{ra:timerbias} sum to at most
\[
 \eta\lambda_-\left(\frac{12+24}{1024}
                         +\frac1{128}+\frac1{128}\right)
 =\frac{13\eta\lambda_-}{256}.
\]
Thus the total is at most
$\eta\mu_u(1/24+1/16+13/256)=119\eta\mu_u/768
<\eta\mu_u/6$.

For the second moment couple to the ideal trial. Away from the censor
band and bad-service event, the output differs from $Z$ by at most
$e=12\Delta/a^2+\omega\le5\eta\lambda_-/256$.
Hence
\[
 \RAE(Y^2\mid\mathcal F)
 \le2\RAE Z^2+2e^2+B^2(8\lambda_+\Delta+\beta).
\]
The first term, averaged over the requested coordinate, is at most
$3N^{-1}\sum_i\lambda_i^2\le3\lambda_+\mu_u$.
The band contribution is at
most $72\eta\lambda_-\lambda_+/1024$.
The lower bound on $a$ gives $B<6\lambda_+\eta^{-1/3}$, so
$B^2\beta\le6\eta^{2/3}\lambda_-\lambda_+/128$.
Together with $\eta\le1/4$ these terms and $2e^2$ are strictly below
$\lambda_+\lambda_-\le\lambda_+\mu_u$.
This proves the bound without replacing the variance
by the much larger worst-case range $B^2$.
\end{proof}

The timer tolerance scales as
$\Delta=\Theta(\eta^{5/3}\lambda_-/\lambda_+^2)$ for a comparable
choice in \eqref{ra:hardware}. Its response and calibration price remain
physical costs. The number of timestamp and arithmetic bits is
$O(\log(1/\eta)+\log\kappa+\log\log(1/\eta))$ beyond the rate and
hardware scale representation. This bit count does not establish that
a device can service that resolution cheaply.

\subsection{A finite confidence bound with dependent repeated cells}
\begin{theorem}[Direct acquired arithmetic-rate certificate]
For each of $K$ declared menu settings, divide the trials into $q$ blocks
of size $m$, where
\begin{equation}
 m=\left\lceil256\kappa\eta^{-2}\right\rceil,
 \qquad
 q=\text{the smallest positive odd integer with }2^{-q}\le\delta/K.
 \label{ra:mom}
\end{equation}
Return the median $\widehat\mu_u$ of the $q$ block averages. Conditional
on the common calibration good event, with probability at least
$1-\delta$, simultaneously for all menu settings,
\begin{equation}
 |\widehat\mu_u-\mu_u|\le\eta\mu_u.
 \label{ra:certificate}
\end{equation}
In fact the proof gives the stronger constant $2\eta/3$ on that event.
The number of trials per menu is
$O(\kappa\eta^{-2}\log(K/\delta))$, independent of $N$ except for
coordinate selection, context access and their actual service costs.
No kinetic rate table, rare stationary count or global trajectory is
acquired by this assay.
\end{theorem}
\begin{proof}
For one block, write $Y_j=m_j+D_j$ with
$m_j=\RAE(Y_j\mid\mathcal F_{j-1})$ and $(D_j)$ martingale
differences. The conditional second-moment bound gives
$\RAE[(\sum_jD_j)^2\mid\mathcal F_{\rm start}]\le4m\lambda_+\mu_u$.
Conditional Chebyshev therefore implies
\[
 \RAP\left(\left|m^{-1}\sum_jD_j\right|>
        \frac{\eta\mu_u}{2}\,\middle|\,\mathcal F_{\rm start}\right)
 \le\frac{16\lambda_+}{m\eta^2\mu_u}
 \le\frac{16\kappa}{m\eta^2}\le\frac1{16}.
\]
Every $m_j$ differs from $\mu_u$ by at most $\eta\mu_u/6$.
Thus each block average is within $2\eta\mu_u/3$ except on an event
whose conditional probability, given all preceding blocks, is at most
$1/16$. For any specified set of $s$ bad blocks, repeated conditioning
bounds their joint probability by $16^{-s}$. A bad median needs at
least $\lceil q/2\rceil$ bad blocks, so its probability is at most
$2^q16^{-\lceil q/2\rceil}\le2^{-q}$. No independence of blocks is
used. Union-bound over the $K$ menus; adaptive scheduling is allowed
provided the menus and trial guarantees are the declared ones.
\end{proof}

\begin{theorem}[A smaller direct-mean schedule]
The same implemented trials can instead be averaged directly. With
\begin{equation}
 n=\left\lceil40\kappa\eta^{-2}\log(2K/\delta)\right\rceil
 \label{ra:bernsteinN}
\end{equation}
trials per menu, their ordinary average satisfies
\eqref{ra:certificate} simultaneously with probability at least
$1-\delta$, conditional on calibration. The bound again holds for
arbitrarily dependent repeated cells satisfying the conditional service
guarantees. No empirical or independence approximation is used.
\end{theorem}
\begin{proof}
For the martingale differences $D_j=Y_j-\RAE(Y_j\mid\mathcal F_{j-1})$,
$|D_j|\le B$ and the conditional variance is at most
$v=4\lambda_+\mu_u$. For $|s|B<3$, expanding the exponential and using
$k!\ge2\cdot3^{k-2}$ for $k\ge2$ gives
\[
 \RAE(e^{sD_j}\mid\mathcal F_{j-1})
 \le\exp\!\left(\frac{s^2v}{2(1-|s|B/3)}\right).
\]
Iterate this conditional bound and apply Chernoff on each side. The
resulting standard martingale Bernstein inequality is
\[
 \RAP\left(\left|n^{-1}\sum_jD_j\right|>t\right)
 \le2\exp\left[-\frac{nt^2}{2(v+Bt/3)}\right].
\]
With $t=\eta\mu_u/2$ and
$B\eta<6\lambda_+\eta^{2/3}\le6\lambda_+$, the exponent is at least
$n\eta^2\mu_u/(40\lambda_+)\ge n\eta^2/(40\kappa)$.
The conditional-mean bias adds at most $\eta\mu_u/6$, giving total
error at most $2\eta\mu_u/3$. Union-bound over the menus.
\end{proof}

If common apparatus or structural calibration fails with probability
$\delta_{\rm cal}$, the unconditional confidence is at least
$1-\delta_{\rm cal}-\delta$. This is distinct from a per-trial conditional
error; reusing a persistently miscalibrated device does not generate new
independent confidence. A common drift or hidden cell memory not covered
by \eqref{ra:trialmoments} invalidates the stated martingale guarantee.

The reported interval for $\mu_u$ is
$[\widehat\mu_u/(1+\eta),\widehat\mu_u/(1-\eta)]$, with outward
rounding; multiply by $N$ for $\Lambda_u$. Choosing the smallest
$\widehat\mu_u$ yields
\[
 \mu_{\widehat u}\le\frac{1+\eta}{1-\eta}\min_u\mu_u.
\]
If the resulting intervals are disjoint, their ordering is certified.
The separate mathematical bridge from boundary escape to the intended
global first-event objective must still be proved at its stated scope;
the local assay alone does not equate stationary flux with lifetime.

The direct-mean implementation is streaming: it stores four current
timestamps, a sum and count, a coordinate index, and the $K$ final
intervals. The median variant additionally stores its $q$ block means.
The arithmetic bit lengths are polynomial in rate and hardware bit
length, $\log N$, $\log n$, and requested precision. Known coordinate
labels and neighborhood access, apparatus cells, context memory and
addressing latency remain priced resources; no free global wiring or
new physical copies are inferred from this small statistical memory.

\subsection{Full dwell and reset bill, including discarded trials}
Let $s=e^{-\lambda D}$ and $q_\lambda=1-s$ for a requested rate.
The ideal sequentially censored trial has expected active dwell
\begin{equation}
 \RAE W_\lambda
 =\frac{q_\lambda}{\lambda}\sum_{j=0}^3q_\lambda^j
 \le\frac4\lambda,
 \qquad W_\lambda\le4D.
 \label{ra:dwell}
\end{equation}
The expected number of pin--release observations is
$\sum_{j=0}^3q_\lambda^j\le4$.
The probability of a discarded censored trial is $1-q_\lambda^4$.
Its entire wasted dwell, including the successful waits before its first
censor, has the exact expectation
\begin{equation}
 \RAE[W_\lambda\RAone_{\rm censored}]
 =\sum_{j=1}^4q_\lambda^{j-1}s
 \left[D+(j-1)\left(\frac1\lambda-
                  \frac{Ds}{q_\lambda}\right)\right]
 \le16D e^{-\lambda D}.
 \label{ra:waste}
\end{equation}
The formulas follow by conditioning on the index of the first censored
wait; $\RAE[T\mid T\le D]=\lambda^{-1}-Ds/q_\lambda$.
Here $q_\lambda>0$ because $D>0$.

With the actual timestamp/deadline service, each normal trial has active
dwell at most $4(D+\Delta)$ and expected active dwell at most
$4/\lambda_-+4\Delta$. A conservative addition
$4\beta(D+\Delta)$ accounts for conditional bad-service branches, whose
timeouts are also hard-capped. Discarded active dwell is bounded by
$16(D+\Delta)e^{-\lambda_-(D-\Delta)}+4\beta(D+\Delta)$ per trial.
All pin, clamp, setup, settling and readout times are additional. These
bounds include unsuccessful exposure; censoring does not make it free.

With $n=mq$ trials for the median version, or \eqref{ra:bernsteinN}
trials for the direct mean, the worst-case active dwell is
$4n(D+\Delta)$ and the expected active dwell is at most
$n[4/\lambda_-+4\Delta+4\beta(D+\Delta)]$.
The complete serial wall-time bound adds, for each trial,
\[
 C_{\rm coordinate}+C_{\rm context}+C_{\rm operating\ setting}
 +4(C_{\rm pin}+C_{\rm settle}+C_{\rm detector}+C_{\rm timer})
 +C_{\rm arithmetic},
\]
plus common calibration, inter-menu finite slew and final deployment.
Use the actual parallel service architecture to convert total work and
component exposure to elapsed time; parallel copies are additional paid
resources. No unpriced simultaneity is assumed. The existing safe mission
preparation and deployment bill are separate from these local target-context
assays and are not removed by a cheaper decision statistic.

The apparatus must detect the \emph{first exit}, not merely read a late
endpoint after possible returns. A failed event record may be flagged
and charged in $\beta$; an unbounded missed-event mechanism cannot be
silently reinterpreted as an exponential wait. If the requested exit
floor is small, its numerical $\lambda_-^{-1}$ exposure is retained.

\section{A matched decision-exposure lower bound at the same local access}
\label{sec:rate_lower}
The quadratic accuracy cost cannot generally be removed by observing more
of a local trajectory, choosing fields adaptively, or granting free
parameter-independent pins. The following argument is an elementary
continuous-time likelihood comparison, in the tradition of sequential
change-of-measure lower bounds. It does not claim that inverse-gamma
estimation or sequential testing is new.

\begin{theorem}[Two-menu decision requires quadratic local exposure]
Fix $0<\eta\le1/4$, $0<\delta<1/2$, $\lambda_0>0$, any $N\ge1$,
and a known $q_*\in(0,1/4]$.
There are two models of $N$ independent bits and two operating menus
such that:
\begin{enumerate}
\item both menus in both models have the same product invariant law
with one-bit probability $q_*$, and the singleton target $b=1^N$ has
$p=q_*^N$;
\item every local \emph{boundary recovery} rate $Q(b,b^{(i)})$ is in
$[\lambda_0(1-2\eta),\lambda_0(1+2\eta)]$;
\item the two models reverse which menu minimizes the mean first-event
time, and also reverse which maximizes it, with a multiplicative gap
larger than $1+\eta$;
\item any adaptive decision procedure with error at most $\delta$ on
both models must have expected total active local exposure
\begin{equation}
 \RAE T_{\rm exposure}\ge
 \frac{(1-2\delta)\log((1-\delta)/\delta)}
      {64\lambda_0\eta^2}
 \label{ra:lower}
\end{equation}
under each model.
\end{enumerate}
The procedure may observe complete continuous local jump paths, use
arbitrary predictable calibrated fields, and select operating menus and
starting bits adaptively. Free resets may reveal their commanded state
but no additional parameter information.
With parallel cells $T_{\rm exposure}$ is the sum of their active
durations, not wall time. The permitted parameter-dependent observations
are the local state/jump records; no extra rate-revealing actuator
side-channel is supplied.
\end{theorem}
\begin{proof}
Put $\theta=2\eta\le1/2$ and
$\lambda_\pm=\lambda_0(1\pm\theta)$.
In model zero, menu one has single-bit $1\to0$ recovery rate $\lambda_+$
and menu two has recovery rate $\lambda_-$. The corresponding $0\to1$
failure rate is $\lambda_\pm q_* /(1-q_*)$; for $q_*=1/4$ it is
$\lambda_\pm/3$. Swap the menus in model one.
Both global generators are common scalar multiples of the same
independent-bit generator. Their equilibrium-safe means therefore scale
exactly inversely with their local flip rate. The ratio is
$(1+\theta)/(1-\theta)>1+\eta$ for every $N$, and the invariant target
mass remains $q_*^N$. Given a strict reliable-family bound
$\bar q<1/2$, choosing $0<q_*\le\min(\bar q/2,1/4)$ places this hard
pair strictly within that family.

To include arbitrary experimental fields, let a field $z$ change the
known conditional probability to $\sigma(\operatorname{logit}(q_*)+z)$,
while the attempt clock in that menu remains $\lambda_\pm/(1-q_*)$.
This is a valid subclass of the
existing access model. The two possible jump intensities at any history
have the form $a=\lambda_0(1+\theta)v/(1-q_*)$ and
$b=\lambda_0(1-\theta)v/(1-q_*)$ in one order, with the same known
$v\in[0,1]$ determined by the field and current bit.
For nonnegative intensities,
\[
 a\log(a/b)-a+b\le (a-b)^2/b
 \le\frac{4\lambda_0\theta^2}{(1-q_*)(1-\theta)}
 \le16\lambda_0\theta^2=64\lambda_0\eta^2.
\]
The same bound applies in the opposite direction; the zero-intensity
limit is interpreted continuously. The controlled point-process
likelihood chain rule integrates this divergence rate over the total
active local exposure. Predictable choices and parameter-independent
resets use the same conditional kernels under the two models and add
no divergence. Consequently
\[
 \RAKL(\mathsf P_0\Vert\mathsf P_1)
 \le64\lambda_0\eta^2\RAE_0T_{\rm exposure},
\]
and conversely with the model indices exchanged. Bounded exposure gives
this directly. For integrable adaptive stopping, stop first at a finite
exposure cap and pass to the limit in the likelihood compensator.

A correct $(1+\eta)$ menu decision must select the appropriate menu in
each model, so it is a binary test with both errors at most $\delta$.
Data processing to that decision gives
\[
 \RAKL(\mathsf P_0\Vert\mathsf P_1)
 \ge \operatorname{kl}(1-\delta,\delta)
 =(1-2\delta)\log((1-\delta)/\delta).
\]
Combine the inequalities in both directions.
\end{proof}

At fixed contrast, the positive assay's expected active exposure is
$O(\lambda_-^{-1}\eta^{-2}\log(K/\delta))$, matching the lower bound's
accuracy and confidence dependence for two menus up to constants and
the stated contrast. The hard worst-case censor cap adds
$\log(1/\eta)$; pins, service latencies and calibration are additional
resources and are not lower-bounded by \eqref{ra:lower}.
The lower bound is uniform in the exponentially small global target
probability. It does not assert optimal dependence on contrast, menu
count, detector price or physical clock resolution.

\section{Ordinary methods, the new interface and retained limits}
The statistic $(n-1)/\sum_{j=1}^nT_j$ for an exponential rate is
classical; the gamma integral above proves the needed $n=4$ instance.
Likelihood-based exponential-rate estimation is an appropriate strong
ordinary comparator. A method that estimates every local rate can also
reach quadratic precision cost per rate, but needlessly acquires $N$
parameters when the decision requires only their average. Conversely,
an ordinary method that groups repeated waits within a randomly chosen
context and then averages the classical unbiased statistics implements
the same direct representation and can match this result. No claim is
made that every conventional estimator has cubic cost.

In particular, a small-time endpoint estimator that divides an observed
exit probability by its hold time uses a time of order $\eta$ to suppress
linearization bias. Its reset/read count can scale as $\eta^{-3}$ even
though its total exposure already scales quadratically. The present
grouped first-exit statistic removes that particular extra reset factor,
retains a finite physical clock construction, and estimates the
arithmetic rate average without a rate table. It does not improve on a
stronger grouped-wait method already using this identity.

The new contribution for the chosen decision is the complete interface:
requested operating contexts, an average of heterogeneous fast-side
rates, finite censor/clip/timing errors including deadline boundary bands,
dependence-safe confidence, all discarded exposure, and a matched
controlled-observation lower bound. Its use in a global rare-event
intervention still relies on the separate target-set and spectral or
potential-theoretic decision theorem. Nothing in the local data certifies
that unobserved failure mechanisms are absent.

\subsection{Composition with the global prevention decision}
In the family of Section~\ref{sec:boundary}, set
$\lambda_-=r_-(1-\bar q)$ and $\lambda_+=r_+$ in the grouped-wait assay.
Let $\widehat u$ minimize its acquired $\widehat\mu_u$. Combining
Theorem~\ref{thm:boundary} with the simultaneous rate intervals gives
\begin{equation}
 \frac{M_{\widehat u}}{\max_uM_u}
 \ge \frac\ell U\frac{1-\eta}{1+\eta}.
 \label{eq:groupedglobal}
\end{equation}
Indeed $M_u=B_u/[N\mu_u p(1-p)]$ with $B_u\in[\ell,U]$,
and the selected mean escape rate is at most $(1+\eta)/(1-\eta)$
times that of every competitor. This is the actual interface between
the local assay and the global objective, not an identification of
stationary flux with lifetime. The confidence is at least
$1-\delta_{\rm cal}-\delta_{\rm dec}$ before the separately reported
preparation abort and deployment allowances. The density and form
allowances below multiply the displayed ratio by
$e^{-2(\beta_{\rm prep}+\zeta_{\rm deploy})}$.

For a prescribed $0<\eps\le1/4$, one explicit allocation is
$\eta=\eps/8$, $\ell/U\ge1-\eps/4$, and
$\beta_{\rm prep}+\zeta_{\rm deploy}\le\eps/8$.
Then each of the three ratio factors is at least $1-\eps/4$,
so their product is at least $1-3\eps/4\ge1-\eps$.
Thus the full deployed decision uses
$O(\kappa\eps^{-2}\log(K/\delta_{\rm dec}))$ grouped trials per
menu after the coarse certificate, with the complete finite service
bill derived above. This accuracy dependence is matched by the local
exposure lower bound at fixed contrast; no optimality is claimed for
the detector, timer or pin prices.
\section{A stronger initialization consequence: the mean from every safe start}
\label{sec:anystart}
The mean-only decision does not require acquiring or preparing an entire
equilibrium initial law in the certified boundary regime. The same finite
certificate controls every safe starting configuration. This changes the
initialization operation while keeping the original mean objective as its
benchmark. It does not supply the earlier marked-law or CDF initial law.

Use the true-generator kernel $P^\circ$, and the $k,L,z,a,g_0,\bar p$
of Section~\ref{sec:boundary}. Define
\begin{align}
 s_*&=(1-\bar p)^2,\qquad
 E_* = La^k+a^{k/2}e^{-g_0L}/g_0,\\
 A_*&=(1-z)(1-\bar p)(1-E_*/s_*),\\
 B_*&=1+(L\bar p+e^{-g_0L}/g_0)/s_*.
 \label{eq:safestartconstants}
\end{align}
Require $E_*<s_*$. All these are certificate arithmetic; $L$ is not an
operating waiting time.

\begin{theorem}[Uniform safe-start mean comparison]
For every menu item and every safe state $x\ne b$,
\begin{equation}
 A_* M_u\le\E_{x,Q^u}\tau_b\le B_* M_u.
 \label{eq:allsafemean}
\end{equation}
The upper bound also holds at $b$, taking $\tau_b=0$. Consequently any
actual starting law $\nu_u$ with $\nu_u(b)\le f$ satisfies
\begin{equation}
 (1-f)A_* M_u\le\E_{\nu_u,Q^u}\tau_b\le B_*M_u.
 \label{eq:imperfectsafemean}
\end{equation}
These are bounds for an unbounded mean, not deductions from a TV estimate.
\end{theorem}
\begin{proof}
For $Z_u(x)=\sum_{j\ge0}[(P^\circ)^j(x,b)-p]$ and $S_u=Z_u(b)$,
\[
 \frac{\E_{x,Q^u}\tau_b}{M_u}
 =(1-p)\left(1-\frac{Z_u(x)}{S_u}\right),\qquad S_u\ge s_*.
\]
At any $x\in C_k$, the finite sum and spectral bounds give
$Z_u(x)\le E_*$. Thus the ratio there is at least
$(1-\bar p)(1-E_*/s_*)$.
For a safe $x$ with $w(x)<k$, the same decreasing superharmonic
function $H$ gives
$\Pr_x(\tau_b<\tau_{C_k})\le H(w(x))\le H(1)\le z$.
Apply the strong Markov property at the first hit of $\{b\}\cup C_k$;
the time already spent is nonnegative. This proves the lower bound.

For every state, $\pi(b)/\pi(x)\le1$. Each of the first $L$ terms of
$Z_u(x)$ is at least $-p$, and the remaining spectral sum has modulus
at most $e^{-g_0L}/g_0$. Hence
$Z_u(x)\ge-Lp-e^{-g_0L}/g_0$, giving the upper bound.
Integrate over $\nu_u$ and use the zero lifetime assigned to an
initially failed state.
\end{proof}

At fixed local margins, choosing the finite tails sufficiently small
makes $A_*/B_*$ approach one as $N$ grows, with the explicit
$z=O(1/N)$ return correction. Therefore this is not an assertion that
all safe starts are equivalent in arbitrary rare-event systems. It is
a consequence of the uniform recovery drift, singleton target and
mixing certificate in this family.

\subsection{The preferred mean-only preparation: a bounded operating wait}
An additional graphical comparison removes inspection when a small
explicit initial-failure mass is allowed by the mean objective. Put
$a_{\rm w}=r_-(1-\bar q)$ and $c_{\rm w}=r_+\bar q$.
Couple the actual failed-bit process $X$ to independent bits $Y_i$,
initially all failed, with recovery rate $a_{\rm w}$ and failure rate
$c_{\rm w}$. A rate-$c_{\rm w}$ clock sets $Y_i$ to failed and gives
$X_i$ its appropriately thinned failure transition. A rate-$a_{\rm w}$
clock recovers both bits, and extra recovery clocks supply the remaining
actual recovery rate. Every actual initial state obeys $X\le Y$, and
these clocks preserve that order. This comparison does not require the
actual interacting system to be attractive; the clocks are a proof
device, not additional physical control.

Thus, from any initial law, at deterministic operating time $t$,
\begin{equation}
 \Pr(X_t=b)\le q_{\rm dom}(t)^N,\qquad
 q_{\rm dom}(t)=\frac{c_{\rm w}}{a_{\rm w}+c_{\rm w}}
 +\frac{a_{\rm w}}{a_{\rm w}+c_{\rm w}}
                  e^{-(a_{\rm w}+c_{\rm w})t}.
 \label{eq:dominatedwarm}
\end{equation}
Every $(N-1)$-coordinate marginal satisfies the analogous bound.
For an actual terminal time in $[D_-,D_+]$, put
$q=q_{\rm dom}(D_-)$ and $\Delta_D=D_+-D_-$.
A visit to $b$ in that window requires either already being there or
an entry caused by one component's failure. Its expected entry count is
at most $Nc_{\rm w}\Delta_D q^{N-1}$. Therefore
\begin{equation}
 \Pr(\exists t\in[D_-,D_+]:X_t=b)
 \le\chi_D:=\min\{1,q^N+Nc_{\rm w}\Delta_Dq^{N-1}\}.
 \label{eq:finitewarm}
\end{equation}
The inclusion of the actual terminal event in this entire-window event
also permits bounded adaptive terminal timing. The actual operating
dynamics must remain in force during commissioning; no field-biased
clock is substituted. The timer must supply a finite upper and lower
elapsed-time bound, and its actual response/calibration cost is paid.

Install and settle the selected intervention, run this finite warm-up,
then begin the mission unconditionally. The actual initial law has
$\nu_u(b)\le\chi_D$, so~\eqref{eq:imperfectsafemean} applies directly.
There is no inspection, rejection, abort or equilibrium-law claim in
this branch. An initially failed mission has lifetime zero. If
$D_-=\log2/(a_{\rm w}+c_{\rm w})$ and the width $\Delta_D$ is fixed,
$\chi_D$ decreases exponentially in $N$, with the explicit linear
prefactor. The physical commissioning time is bounded independently
of $N$ at fixed local margins, while summed component exposure is
$ND_+$; installation, propagation, finite slew, settling and timer
services remain additional. Repairability during commissioning remains
essential. If the declared target loss cannot accommodate $\chi_D$,
use the next interface or the stronger paid relative preparation.

\subsection{A finite one-bit alternative at the actual operating setting}
This physical realization requires a repairable installation and a
noninvasive meter for one chosen actual component. The earlier full
safe-state inspection access supplies such a meter; local copies alone
would not. The meter must capture that component's binary state at a
declared sample time and report by a finite deadline. A finite capture-time
uncertainty is included in the elapsed-delay allowance below. This is
stated access, not a newly invented sensor. Its response, calibration,
latency, setup and read-error prices remain explicit. The capture time and
its finite jitter are exogenous to the current focal path conditional on
previous history; a path-adaptive capture needs an additional error bound.

Let the system evolve at the selected actual menu from any starting
configuration. For the chosen bit, put
\[
 a_{\rm w}=r_-(1-\bar q),\qquad c_{\rm w}=r_+\bar q.
\]
Its instantaneous failed-to-working rate is at least $a_{\rm w}$ and
its working-to-failed rate is at most $c_{\rm w}$, for every neighborhood.
If $v(t)$ is its working probability, then
\[
 v'(t)\ge a_{\rm w}(1-v(t))-c_{\rm w}v(t).
\]
No independence of the components is used. Thus an exogenous warm-up
of duration at least $t_0$ gives
\begin{equation}
 v(t)\ge v_0:=\frac{a_{\rm w}}{a_{\rm w}+c_{\rm w}}
       \left(1-e^{-(a_{\rm w}+c_{\rm w})t_0}\right).
 \label{eq:witnesswarm}
\end{equation}
The bound is uniform over the entire history at the start of an attempt.
A convenient finite choice is
$t_0=\log2/(a_{\rm w}+c_{\rm w})$, with an outward-rounded timer.

Suppose the captured-bit error is at most $\beta_{\rm read}$ conditional
on its true state and previous history, and the mission begins no later
than $\Delta_{\rm read}$ after the earliest possible capture time. The
operating dynamics continue during this interval. Accept only a reported
working bit. Then
\begin{equation}
 s_{\rm acc}\ge(1-\beta_{\rm read})v_0,\qquad
 f_{\rm start}\le
 \frac{\beta_{\rm read}+c_{\rm w}\Delta_{\rm read}}
      {(1-\beta_{\rm read})v_0}.
 \label{eq:witnesserrors}
\end{equation}
For the second inequality, an accepted but globally failed terminal
state requires either a wrong captured-bit report or a subsequent failure
of that bit. A correct working bit has probability at most
$c_{\rm w}\Delta_{\rm read}$ of such a transition. Dividing by the
acceptance lower bound proves the claim. An erroneous acceptance is
therefore retained as possible initial failure, not relabeled a certainly
safe preparation. Detector common-mode failures must be charged in the
common calibration confidence or the stated conditional bound.

Repeat up to
\[
 J_{\rm prep}=\left\lceil
 \frac{\log\delta_{\rm abort}}{\log(1-s_{\rm acc})}
 \right\rceil
\]
attempts; return an explicit abort if none accepts. Uniform conditional
bounds give abort probability at most $\delta_{\rm abort}$ even for
dependent reused apparatus. The mixture over accepted attempts still
has $\nu_u(b)\le f_{\rm start}$. No equilibrium sampler, rare-state
rejection, complete state readout or full-system pin is required.

Each attempt pays its actual whole-system warm-up and latency time,
one capture/read, controller/timer response, and any reset/setup.
Summed component exposure is at most
$N J_{\rm prep}(t_{\rm warm}^{+}+\Delta_{\rm read})$, with other active
installation intervals added. At fixed local bounds and fixed desired
$f_{\rm start},\delta_{\rm abort}$ this is $O(N)$; it is not constant
total exposure merely because the installation evolves in parallel.
Field installation, propagation, slew, settling and the stated meter's
physical price are additional. If the source lacks the required one-bit
meter, or its latency/read fidelity cannot satisfy~\eqref{eq:witnesserrors},
this realization is unavailable and the separately paid preparation must
be used. The process must be repairable before mission start.

\subsection{The complete deployed decision factor}
Let $F_{\rm dec}$ denote the proved selected-to-best ratio for the
equilibrium-safe benchmark, for example
$(\ell/U)(1-\eta)/(1+\eta)$ in~\eqref{eq:groupedglobal}.
For the just-constructed preparations, the actual means
$\mathcal M_u=\E_{\nu_u,Q^u}\tau_b$ obey
\begin{equation}
 \boxed{\frac{\mathcal M_{\widehat u}}{\max_u\mathcal M_u}
 \ge (1-f_{\rm start})\frac{A_*}{B_*}F_{\rm dec}.}
 \label{eq:deployedsafe}
\end{equation}
For the unconditional warm-up branch take $f_{\rm start}=\chi_D$ and
there is no preparation abort. For the one-bit branch and a mission
definition assigning lifetime zero to a preparation abort,
multiply the lower guarantee by $1-\delta_{\rm abort}$. The confidence
of the acquisition/decision certificate remains a separate probability.
All displayed factors are checked against the requested tolerance before
issuing the decision. Additional deployed-generator error is allowed only
under its separately proved common-law form bounds or by widening the
actual operating certificates; it is not hidden inside detector error.

This completes the mean-only initialization interface. Relative equilibrium
preparation remains a distinct stronger operation for outputs that require
that initial law. The next section retains its conservative bill; the
independent logarithmic-Sobolev construction in the archive improves its
mixing schedule under its additional acquired norm condition.

\section{Acquisition, repeated cells, preparation and the full service bill}
\label{sec:acquisition}
The decision uses coarse certificates, not a finely calibrated model. Their acquisition remains part of an end-to-end claim. Already valid engineering or previously acquired certificates can be reused, but their provenance and costs are not erased.

\subsection{Coarse probabilities and actual operating clocks from paired holds}
At one actual menu environment and neighbor context, let $r$ and $q$ be the unknown two-state clock and equilibrium failure probability. Pin to zero and read after a hold to obtain a $0\to1$ indicator; separately pin to one for a $1\to0$ indicator. At a common exact hold $t$ their means are
\begin{equation}
 a=q(1-e^{-rt}),\qquad b=(1-q)(1-e^{-rt}).
 \label{eq:calpair}
\end{equation}
Neither experiment waits for a locally rare natural starting bit. The two paid pins provide it. Their sum exposes the refresh clock without a Boltzmann suppression, and their ratio exposes $q$.

The hold lengths need not be identical. Suppose each is exogenous to its current focal path and belongs to $[t_-,t_+]$, with $t_->0$. They may depend on previous experimental history. Let simultaneous confidence intervals for the conditional average indicator means be $[a_-,a_+]$ and $[b_-,b_+]$, including all pin/read/context mean biases. Put $c_t=t_+/t_-$. Concavity of $1-e^{-rt}$ implies that the ratio of any two admissible hold responses is between $c_t^{-1}$ and $c_t$. Thus valid intervals include
\begin{align}
 q&\in\left[\frac{a_-}{a_-+c_t b_+},\,
             \frac{c_t a_+}{c_t a_++b_-}\right],\label{eq:qinterval}\\
 r&\in\left[\frac{-\log(1-d_-^{\rm obs})}{t_+},\,
             \frac{-\log(1-d_+^{\rm obs})}{t_-}\right],
 \quad d_-^{\rm obs}=\max(0,a_-+b_-),\quad
 d_+^{\rm obs}=\min(1,a_++b_+).
 \label{eq:rinterval}
\end{align}
A zero denominator gives the conservative endpoint zero or one; $d_+^{\rm obs}=1$ gives an infinite upper clock bound. A supplied physical upper rate bound remains useful. No finite observation can certify an arbitrary unbounded rate above a saturated detector's range.

For $m$ observations of either starting bit, martingale Hoeffding gives statistical radius $\sqrt{\log(2C/\delta_{
\rm cal})/(2m)}$ around its conditional average; add the certified conditional mean bias. This is valid for reused cells with the declared memoryless conditional response. It does not infer memorylessness from a short record.

Use one menu item's probability intervals to certify the common equilibrium law. Enumerate the declared local contexts. For neighboring contexts $s,s'$ differing at $j$, form the maximum difference allowed by their $q$ intervals. Maximize over the remaining local bits to obtain $\overline c_{ij}$. Verify
\begin{equation}
 \max_{i,s}q_i^+(s)\le\bar q<1/2,\qquad
 \max_i\sum_j\overline c_{ij}\le1-\alpha.
 \label{eq:coarsecertificate}
\end{equation}
Clock lower bounds must hold in every declared context of every menu item. Their enumeration costs $C_{
\rm ctx}=K\sum_i2^{\deg(i)}$, not its logarithm. Common equilibrium across items is a structural premise, not a consequence of similar fitted probability intervals. These bounds give $\gamma=\alpha r_-$.

If the certified probability intervals also give $q_i(s)\ge q_->0$, then every configuration has $\pi(x)\ge q_-^N$, since $1-q_i(s)\ge1-\bar q>q_-$ after reducing $q_-$ if necessary. This supplies the mass floor for relative preparation. Otherwise the accepted bounded-energy learner can supply that floor at its separately paid field/dwell cost. The decision itself uses no lower bound on $q_i$.

Under fixed strict margins, all intervals above can be obtained with constant per-context probability accuracy. If $\Delta_{
\rm cal}$ denotes a sufficient additive probability accuracy, the coarse acquisition uses
\begin{equation}
 O\left(C_{\rm ctx}\Delta_{\rm cal}^{-2}
            \log\frac{C_{\rm ctx}}{\delta_{\rm cal}}\right)
 \label{eq:calcost}
\end{equation}
pins and binary reads, and at most $t_+$ times that count in free-cell exposure. The actual hold choice and all context, timer and field costs are added. A sequence of decreasing accuracies with summable confidence can discover strict slack; an exact boundary need not terminate. A promised positive slack gives a finite stated schedule. In particular, no global $1/p$ factor is introduced, but rare local responses, large degree or small margins can make calibration expensive.

\subsection{What a local service costs}
\label{sec:services}
For the preferred boundary branch, a single call costs
\begin{equation}
 C_{\rm setup}(d,u)+C_{\rm pin}(\beta)+C_{\rm field}(\xi_u)
 +C_{\rm slew}+h(1+\omega)+C_{\rm read}(\beta)+C_{\rm timer}(h\omega)
 +C_{\rm bits}.
 \label{eq:physicalservice}
\end{equation}
These symbols denote actual validated services, not unit-cost promises. Context pins are paid, including setting the neighbors to the failed value. The system itself is never prepared globally failed: only one focal local copy and its declared neighborhood are set. The operating environment is reproduced at every experiment and on deployment. If a small field error can change an unknown clock discontinuously, field precision alone gives no clock guarantee; a reproduced environment or a proved operating-rate envelope is required.

Boundary trials need only conditional probability bias bounded by the decision-level allowance in Section~\ref{sec:boundary}. It does not shrink as the number of trials grows. Dependence from repeated cells is handled by the conditional concentration theorem. The exogenous timer may drift within its declared range as a function of past observations. A separate weaker bound covers a stopping time responsive to the current focal path. Neither permits unbounded detector delay or missed exits without its stated allowance.

The return branch requires a fixed reversible ideal $P_h$. A finite cap bounds its total calls by $D_{
\rm calls}=KnJ$. Conditional pin, read, numerical and selector errors with sum at most $\delta_{
\rm service}$ couple it to that ideal experiment except with that probability. This branch therefore may need finer service than the direct boundary branch. The distinction is intentional. Finite random-coordinate selection is implemented by binary rejection with a cap; its bias can be charged at the mean level for the boundary branch, or in the total coupling ledger for returns. Integer counters, rational interval arithmetic and output confidence bits are paid.

There is no universal claim that $C_{\rm pin}(\beta)$ or $C_{\rm read}(\beta)$ grows logarithmically. If a given instrument supports a stable captured bit and conditionally bounded independent read faults below one half, repeated reads and majority give that familiar logarithmic amplification. The capture and fault assumptions are additional supplied hardware in that realization. The theorem otherwise keeps the actual service function.

\subsection{The separate stronger relative-equilibrium initialization}
The mean-only boundary decision can instead use Section~\ref{sec:anystart}.
This subsection retains the full relative-law preparation for the return
fallback and the older forecasts that genuinely require it.
For a selected menu item let the repairable operating system mix, inspect it, retain a safe state and otherwise retry. The process must remain available through failed states during preparation; this is not free preparation of an irreversibly destroyed specimen. The stationary safe mass is at least $1-\bar p\ge1/2$.

With $\pi_{
\min}\ge q_-^N$, reversible spectral mixing gives, from any initial state,
\[
 \left|\frac{P_t(x,y)}{\pi(y)}-1\right|\le e^{-\gamma t}/\pi_{
\min}.
\]
Choose $t_{
\rm prep}\ge\gamma^{-1}\log[1/(\xi\pi_{
\min})]$ and set $\xi\le\tanh(\beta/2)$. Conditioning a correctly inspected safe state then yields
\begin{equation}
 e^{-\beta}\le\frac{d\nu_u}{d\pi(\cdot\mid A)}\le e^\beta.
 \label{eq:prep}
\end{equation}
A finite $O(\log(1/\delta_{
\rm prep}))$ retry cap suffices at the stated safe mass, with explicit preparation abort. Inspection and hold errors must be included using the accepted finite relative-preparation construction. A smaller safe-mass floor would require its inverse in the retry cost. Here the locally certified $\bar q$ supplies the safe-mass slack directly.

This pointwise density bound, not just TV, is required for an unbounded mean:
$e^{-\beta}M_u\le\E_{
\nu_u}\tau_b\le e^\beta M_u$.
If the actual deployed generator additionally has a common-invariant form discrepancy $e^{\pm\zeta}$ from the calibrated one, the ratio guarantee loses at most $e^{-2(\beta+\zeta)}$. A mean-zero preparation-abort outcome multiplies an unconditional lower guarantee by its success probability; for a conditional-on-success mission, the abort probability is reported separately. General changes to the invariant law do not enjoy this cancellation and require the earlier acquired uncertainty machinery.

Only the selected action is physically prepared. Acquisition precedes the mission. Protecting a live installation during acquisition needs its own priced hold, shutdown or exposure analysis. The intervention is not instantaneous feedback: computation, installation, field travel, finite slew, settling and inspection all take time. The deployment bill is
\begin{equation}
 C_{
\rm install}+C_{
\rm ramp}+C_{
\rm settle}
 +O\left(\log(1/\delta_{
\rm prep})\,[t_{
\rm prep}+C_{
\rm inspect}(N)]\right).
 \label{eq:deployment}
\end{equation}
Under this displayed sufficient relative-mixing bound, the free-evolution wall time is $\widetilde O(N)$ at fixed margins. It acts on an entire $N$-component installation and therefore uses $\widetilde O(N^2)$ summed component exposure. The inspection and other service functions are additional. It is not a lower bound on possible preparation time. The acquired-information separation does not justify silently dropping this deployment cost.

\subsection{Resource scaling and nonclaims}
At fixed degree, clock contrast, reliability margin, decision accuracy and confidence, coarse calibration has $\widetilde O(KN2^d)$ local calls. The boundary decision has a number of local calls independent of $N$; the return branch has $O(K\operatorname{polylog}N)$, with the explicit accuracy and numerical margin factors in its schedule. A sparse set of recovered indices represents a return configuration; initializing $b$ requires an empty set, not $N$ physical failures. The known interaction graph and full menu description still occupy their actual input size. Loading them can cost $\Omega(N)$ even when a subsequent decision samples only a few coordinates.

These are local experiment and bit-work statements, not claims of constant wall time for a growing apparatus. Parallel acquisition pays the number of cells, controllers and detectors; the prediction lower bound charges their summed exposure. Large sufficient calibration constants can dominate the short decision. The numerical executions below do not establish economical general hardware.

\section{What this earns beyond an absolute-forecast pipeline}
The consequence is a decision-specific acquisition theorem. Under its stated structural class, an intervention can be selected with a relative first-event mean guarantee while the absolute lifetime remains statistically unresolved. The missing scale is not approximated and then forgotten: it cancels from the exact comparison. The boundary branch also avoids constructing any experimental Markov chain. Its finite local binary responses are mapped directly into a decision interval.

\subsection{The strongest applicable ordinary methods}
Green identities, Dirichlet and Thomson principles, capacity-based rare-event analysis, and exponential rare-hit approximations are established \cite{aldousfill,bovier,aldousbrown}. The boundary-flux method is an especially strong ordinary comparator in the uniformly reliable large-$N$ family. It is the preferred branch here. It would be misleading to compare only with brute-force waiting for a catastrophe or with reconstructing $2^N$ states.

Loper, Zhou and Geman already reduce suitable rare first-passage probabilities to target-neighborhood capacities and local simulations, with explicit flatness errors \cite{loper}. Their paper also explains why those capacities do not recover an overall first-passage time scale. The present operation exploits a common invariant law so that this missing scale cancels in an intervention comparison, proves a finite Hamming-neighborhood correction from local bounds, and constructs the measurements using unknown operating local cells. We do not claim that local capacity methods or nuisance cancellation are new ideas in themselves.

For independent identical components an ordinary birth--death or product-integral solver is computationally superior to a general rare-path engine. For known numerical coefficients it can give an accurate absolute mean quickly. The lower bound below concerns acquiring its unknown coefficients sufficiently well: granting the best such exact solver does not supply the needed information. It compares with every absolute-mean acquisition procedure in the same local access class, not just a naive uniform confidence-box estimator.

The accepted marked sampler and CDF certificate remain stronger output capabilities. They support general event predicates and other losses under their scopes. This new operation saves work by committing to one objective, one common-invariant intervention family, and a singleton event; it does not make those broader forecasts unnecessary for questions they actually answer. No universal superiority over every specialist direct-decision algorithm is claimed. The explicit boundary branch is one such specialist method, now joined to finite acquisition, errors and deployment.

\section{A same-access separation between a decision and an absolute forecast}
\label{sec:sameaccesslower}

This lower bound concerns acquired information, not the running time of a
solver given exact weights.  It grants the forecaster strictly more
observations than the finite-hold decision needs.  In particular it grants
perfect clocks, arbitrary calibrated additive fields, and complete continuous
observation of each released local cell.  The obstruction therefore does not
come from a slow unknown clock or from requiring an inefficient parameter
estimator.

\paragraph{The experiment class.}
There are $n$ independent binary components, with unknown logits
$\theta_i$, equilibrium failure probabilities $q_i=\sigma(\theta_i)$,
and \emph{known unit refresh clocks}.  At applied field $u$, a local cell
has rates
\[
  0\longrightarrow1:\ \sigma(\theta_i+u),\qquad
  1\longrightarrow0:\ 1-\sigma(\theta_i+u).
\]
An experiment may adaptively choose a component, an arbitrary field, a
known initial bit, a release duration, and its stopping rule.  It may observe
the complete resulting local path.  Several cells may operate in parallel;
their free durations are added.  A supplied fresh equilibrium bit at a
chosen field is also permitted and counted as one local observation.  Pins,
controller randomization, and computations convey no additional unknown
parameter by themselves.  An equilibrium bank of $n$ bits counts as $n$
local observations; parallelism does not turn its $n$ sensing channels into
one unit of acquired evidence.  These are the declared local state/clock
experiments, not direct measurements of unknown Hamiltonian coefficients.

Write
\[
  W=B+r^*\sum_{\text{released cells}} \text{free duration},
\]
where $B$ is the number of supplied fresh equilibrium bits and $r^*$ is
the known maximum available refresh clock (one in the unit-clock model).
Allowing additional known clock modules up to a fixed $r^*$ only changes
this normalization: their rates are the same under the two hypotheses.
This relaxed
bill omits pin, detector, field, memory, and computational costs, so a lower
bound on it remains a lower bound when those services are charged.

\paragraph{Theorem (forecast information lower bound).}
Let $n\ge4$, $0<\epsilon\le1/8$, and $0<\delta<1/2$.
Suppose an adaptive procedure returns an estimate $\widehat M$ satisfying
\[
 \Pr_\theta\{|\log\widehat M-\log M(\theta)|\le\epsilon\}\ge1-\delta
\]
for every allowed model, where $M$ is the mean first visit to
$b=(1,\ldots,1)$ from the equilibrium law conditioned on $b^c$.
Then, on at least one allowed model,
\begin{equation}
  \mathbb E W\ \ge
  \frac{n^2}{512e\,\epsilon^2}
       \operatorname{kl}(1-\delta,\delta).
  \label{eq:forecastlower}
\end{equation}
With only fresh equilibrium-bit experiments the stronger constant
$n^2\operatorname{kl}(1-\delta,\delta)/(128\epsilon^2)$ holds.
For fixed $\delta<1/2$, this is an $\Omega(n^2/\epsilon^2)$ local-evidence
requirement.  It holds even though the operating gap equals one, the clocks
are known, and every coordinate can be actively balanced.

\paragraph{Proof.}
First consider equal logits $\theta_i=\theta$, and put $q=\sigma(\theta)$,
$p=q^n$.  Independence and the two-state semigroup give
\[
 P_t(b,b)=\{q+(1-q)e^{-t}\}^n.
\]
The classical fundamental-matrix identity, also derived in the return
certificate below, gives
\begin{align}
 S_n(q)&=\int_0^\infty[P_t(b,b)-p]\,dt
       =\sum_{k=1}^n {n\choose k}
                    \frac{q^{n-k}(1-q)^k}{k},\\
 M_n(q)&=\frac{S_n(q)}{p(1-p)}.
 \label{eq:productmeanlower}
\end{align}
Let $K\sim\operatorname{Bin}(n,1-q)$.  Then
$S_n=\mathbb E[K^{-1};K>0]$.  Since
$1/k\le2/(k+1)$ for $k\ge1$,
\[
 S_n(q)\le\frac{2}{(n+1)(1-q)}.
\]
Direct differentiation of the finite sum gives
\begin{equation}
 -\frac{d}{d\theta}\log M_n(\sigma(\theta))
 =\frac{1-p}{S_n(q)}-\frac{np(1-q)}{1-p}.
 \label{eq:meanlogderivative}
\end{equation}
For $n\ge4$ and $q\le1/2$, the right side is at least
\[
 \frac{15(n+1)}{64}-\frac n{15}\ \ge\ \frac n8.
\]
Take $\theta^{(0)}_i=-\log3$ and
$\theta^{(1)}_i=-\log3+d$, where $d=32\epsilon/n\le1$.
The probabilities along this segment stay below $1/2$.
Equation~\eqref{eq:meanlogderivative} therefore yields
\[
 \log M(\theta^{(0)})-\log M(\theta^{(1)})\ge4\epsilon.
\]
A valid estimate distinguishes these two hypotheses with both error
probabilities at most $\delta$.

For one equilibrium-bit observation at any adaptively selected field,
the Bernoulli log-partition function has second derivative at most $1/4$.
Consequently its conditional divergence between the two hypotheses is at
most $d^2/8$, independently of the field magnitude.  For a released cell,
write $\lambda_0(x,u),\lambda_1(x,u)$ for the applicable exit rates.
Because the derivative of each logistic log rate has magnitude at most one,
\[
 \left|\log\frac{\lambda_1(x,u)}{\lambda_0(x,u)}\right|\le d.
\]
The continuous-path likelihood identity gives conditional divergence
\[
 \mathbb E_0\int
 \left[\lambda_0\log\frac{\lambda_0}{\lambda_1}
             -\lambda_0+\lambda_1\right]dt
 \le\frac{ed^2}{2}\,\mathbb E_0\int dt,
\]
since $\lambda_0\le1$ in the unit-clock case (or at most $r^*$ with
additional clock modules) and
$e^v-1-v\le e d^2/2$ for $|v|\le d\le1$.
Adaptive fields, adaptive stopping, and repeated use of a cell obey the
same conditional likelihood chain rule.  Truncating an unbounded stopping
rule and then taking a limit gives
\[
 \operatorname{KL}(\mathsf P_0^{\rm data}\Vert\mathsf P_1^{\rm data})
 \le\frac{ed^2}{2}\,\mathbb E_0 W.
\]
Data processing through the final binary test lower-bounds this divergence
by $\operatorname{kl}(1-\delta,\delta)$, proving
\eqref{eq:forecastlower}.  The equilibrium-only bound uses $d^2/8$ in
place of $ed^2/2$.  Finite instrument errors cannot increase the information
available beyond this deliberately idealized observation class.
\hfill$\square$

\paragraph{Heterogeneous and nontrivial decision instances.}
The hard pair need not consist of identical components.  Let initial
$q_i$ belong, for example, to $\{1/4,1/3\}$ in any unknown arrangement.
Apply the same additive logit shift $t$ to every component.
If $K$ is the resulting Poisson-binomial number of working components,
$S=\mathbb E[K^{-1};K>0]$ and $p=\prod_iq_i$ satisfy
\[
 \frac{dS}{dt}=\Lambda S-(1-p),\qquad
 \frac{d\log p}{dt}=\Lambda,
 \qquad \Lambda=\sum_i(1-q_i).
\]
Thus
\[
 -\frac{d\log M}{dt}=\frac{1-p}{S}-\frac{p\Lambda}{1-p}.
\]
As long as all $q_i\le1/2$, stochastic domination by
$\operatorname{Bin}(n,1/2)$ gives the same lower bound $n/8$.
For $n\ge8$, $\epsilon\le1/8$, the shift $32\epsilon/n$ keeps the
displayed heterogeneous family below $1/2$.  The information bound is
unchanged.  These models can have a nontrivial optimal assignment of unequal
clock modules, while the unknown common energy shift leaves their optimal
ordering unchanged.  The lower bound is therefore not an artifact of
uncertainty about a globally rescaled time unit.

\paragraph{Matching order for prediction in this product family.}
The $n^2$ dependence is attainable up to confidence and finite-service
logarithms; it is not an inverse-$p$ assertion.  For $q_i$ in a fixed compact
subinterval of $(0,1)$, sample each local equilibrium probability $m$ times,
clip its empirical frequency to a slightly larger compact subinterval,
and add the $n$ log frequencies.  The bias is $O(n/m)$ and bounded
differences give fluctuation
$O(\sqrt{n\log(1/\delta)/m})$.  Taking
$m=O(n\epsilon^{-2}\log(1/\delta))$ estimates $\log p$ to $O(\epsilon)$
using $O(n^2\epsilon^{-2}\log(1/\delta))$ local observations.  The
nonrare return prefactor $S$ is acquired separately by the finite return
operation.  Every ideal equilibrium read can be realized by a finite hold
at the known unit clock, with conditional bias below its allotted error;
this introduces an explicit logarithmic dwell factor, not $1/p$.

\paragraph{Meaning of the separation.}
At fixed margins and fixed decision accuracy, the direct acquired
intervention constructions use $\widetilde O(n)$ local evidence, while
\emph{every} uniformly accurate absolute-mean forecaster in the same local
access class has the $\Omega(n^2)$ requirement above.  A method is free to
use the same nuisance cancellation and become a direct decision method;
the theorem does not claim a lower bound against such a method.  It shows
that requiring a calibrated absolute forecast before deciding is an
additional statistical task.  It is not a claim that ordinary simulation
is the strongest baseline, nor that all rare-event decisions are cheap.

The proof uses classical likelihood/change-of-measure tools, as in
fixed-confidence experimental identification.  The established
fundamental-matrix and bandit results are machinery; no priority claim is
made for those identities.  The derived separation is for the explicit
rare first-event intervention family and the paid local experiment class
above.


For direct comparison with the numerical family's $\bar q=1/4$, the same proof can use $\theta^{(0)}=-\log7$ and $N\ge8$. Its shifted probabilities remain below $1/4$ for $\epsilon\le1/8$, while the derivative and information bounds used above still hold. Thus the quadratic forecast lower bound also holds inside a fixed uniformly reliable subclass with the same kinds of coarse certificates, known fast clocks and safe preparation. The output bit length, input description and parallel hardware costs are not removed by either method.

\section{Executed growing-family decision and its limits}
\label{sec:execution}
The root program implements only the data flow required by the decision. Its source simulator is hidden behind an actual-context local-hold interface. It supplies binary outcomes, and the decision receives public coarse bounds, menu labels and $N$. It never reads the numerical energies, clocks, stationary mass or absolute lifetime.

The family has nonuniform fields $h_i=-2+0.25\sin(0.73i)$ and ring interactions $J_{i,i+1}=0.12\cos(0.41i+0.3)$. The three clock menus cross between contexts. One has clocks in $[0.25,0.30]$ when both neighbors have failed and clocks in $[0.85,1]$ elsewhere; a second reverses those assignments. A third has a different heterogeneous context response in $[0.5,0.65]$. Consequently there is no coordinatewise generator order deciding the menu. All share the same Gibbs law. Valid loose certificates are $\bar q=1/4$, $r_-=0.2$, $r_+=1$, $\alpha=0.7$ and $\gamma=0.14$.

Each run uses $4096$ binary local holds for each of three actions. The nominal hold is $h=1/4$, its exogenous deadline varies with prior history by at most $0.0025$ time units, and the simulator adds read corruption with probability $0.0005$. The certificate permits combined conditional bias $0.001$. It uses an exact-integer Chernoff inversion and exact-rational structural inequalities. A tail budget $1/1200$ per item, tail and size gives simultaneous statistical confidence at least $0.98$ across all four runs in the ideal random experiment. The executed seeded PRNG record is a reproducible diagnostic, not a certified physical random source.

\begin{center}
\begin{tabular}{rrrrr}
\toprule
$N$&Boundary $k$&Upper $q_0S$&Mean ratio vs fast menu&Mean ratio vs mixed menu\\
\midrule
64&14&1.053914&1.735710&1.045747\\
256&15&1.009896&1.860830&1.055450\\
1024&17&1.002926&1.904448&1.257433\\
4096&18&1.002223&1.717734&1.080755\\
\bottomrule
\end{tabular}
\end{center}
The last two columns are conservative certificate lower bounds on the selected action's true equilibrium-safe mean divided by the competitor's mean, before the separately specified deployment/preparation allowance. They are not estimates of the actual improvements. Every run certifies the same menu as best. The state spaces range from $2^{64}$ to $2^{4096}$; the proof, rather than enumeration, governs the growing family.

Every size uses $12\,288$ local holds and $3061.76$ total free-cell time units. No full-system failed preparation and no virtual trajectory were used. Additional pin, neighbor-context, read, timer, field, setup, controller and installation services are not included as zero. The experiment assumes the validated coarse certificates; their full finite acquisition is constructed above and its conservative bill is evaluated separately below.

\subsection{Composition with the stronger mean-only initialization}
The file \texttt{deployment\_composition.json} applies the all-safe-start
theorem to the same recorded counts using exact rational arithmetic.
The preferred branch runs the actual installation for a time in
$[2,2.001]$ with no inspection or rejection. The graphical comparison
gives initial-failure mass at most
$(13/16)^N+(N/4000)(13/16)^{N-1}$, including the entire clock window.
It is below $1.727\,10^{-6}$ at $N=64$ and decreases exponentially.
The actual selected-to-best mean is then at least $0.9904350$ at
$N=64$; the chosen menu is certified best at the other three sizes.
Whole-system free evolution is capped at $2.001$ model-time units,
and summed component exposure at $2.001N$. This evaluated interface
uses the proved uniform rate bounds and repairable commissioning;
it was not executed on physical hardware and does not assert a
certainly safe initial specimen or equilibrium-law preparation.

A separately evaluated meter contract has captured-bit error at most
$1/4000$, capture-to-mission delay at most $0.001$, and actual warm-up
in $[2,2.001]$ model-time units. It gives working probability at least
$3/16$, acceptance at least $11997/64000$, and conditional initial
failure at most $32/11997<0.002668$. Thirty-four attempts bound the
preparation-abort probability by $0.001$. Their whole-system warm-up
and latency cap is $68.068$ time units, costing $68.068N$ component-time
before the other physical services. This is an evaluated interface,
not executed physical preparation.

The selected-to-best actual successful-preparation mean is certified
at least $0.9877949$ at $N=64$. At $N=256,1024,4096$, the selected
menu is still certified best after these additional allowances. If
preparation abort receives zero lifetime, the corresponding selected-to-best
guarantees are $0.9868071$ at $N=64$ and at least $0.999$ for the other
sizes. Thus the equilibrium-safe table is not silently reused as an
unchanged physical-start certificate. No full initial law is claimed.

\subsection{Separate grouped-wait execution}
The completed \texttt{next6\_rate\_assay} program executes finite
censoring, clipping, persistent clock offsets and exact rational output
arithmetic on heterogeneous local contexts at $N=64,4096,262144$.
It uses $49\,152$ groups per menu at relative tolerance $1/8$, for six
menu assays with simultaneous confidence at least $0.95$ conditional
on the coarse bounds. All three pairs have disjoint confidence intervals.
The six executions used $294\,912$ groups, $1\,179\,647$ pin--release
observations and $1\,597\,822.63$ active model-time units. The one
censored group and its discarded exposure are included. Observed relative
errors below $0.54\%$ do not replace the $12.5\%$ certificate.
This is a local-interface execution, distinct from the interacting-ring
global decision and from a physical apparatus experiment.

\subsection{A full bill can still be large}
The file \texttt{calibration\_costs.json} evaluates a simple sufficient coarse acquisition with hold $5$, timing error at most $0.05$, per-indicator statistical radius $0.01$, service bias $0.0005$, and total calibration error probability $0.005$. It includes all $4N$ contexts for each of three clock menus and both starting bits. This schedule was evaluated, not experimentally executed.
\begin{center}
\begin{tabular}{rrrr}
\toprule
$N$&Coarse pins/reads&Coarse free exposure&System prep time cap\\
\midrule
64&$1.02364\,10^8$&$5.16936\,10^8$&$1.89661\,10^4$\\
256&$4.52039\,10^8$&$2.28280\,10^9$&$7.46155\,10^4$\\
1024&$1.97849\,10^9$&$9.99138\,10^9$&$2.97213\,10^5$\\
4096&$8.59541\,10^9$&$4.34068\,10^{10}$&$1.18760\,10^6$\\
\bottomrule
\end{tabular}
\end{center}
Exposure is in model time units, not automatically seconds. The preparation column is whole-system free-evolution time, using a pointwise density-log allowance $0.01$, a certified conditional floor $q_-=0.025$, and eleven capped attempts. Its summed component exposure is $N$ times that column; inspection, installation and field ramping are additional. If one time unit were one second, the first coarse exposure alone would exceed sixteen serial cell-years. This is a deliberately exposed sufficient cost, not a lower bound and not a claimed useful device. The inexpensive postcertificate decision must not be presented as the entire acquisition bill. The theorem's established advantage is asymptotic information scaling and removal of unnecessary forecast requirements; a practical regime still depends on the actual source and services.

\subsection{Useful unresolved attempts retained}
The first numerical certificate used the larger class of timers allowed to stop in response to the current focal path. Its weaker coefficient left three of four menu comparisons unresolved at the chosen sample count. That record is preserved in \texttt{stronger\_timer\_contract\_unresolved.json}. The executed source uses exogenous deadlines; applying the appropriate, separately proved bound certifies all four using the same underlying outcome counts. The final coarse envelopes were also loosened to have explicit acquisition slack. No outcome was discarded and no numerical failure was renamed a guarantee.

The initial candidate architecture used return probes for every decision. The ordinary boundary-flux comparator made that unnecessary in the certified regime and replaced it. Earlier attempts to use only generator order or convexity were retained as limited mechanisms: they do not solve arbitrary preventive allocation. The distinctions between a mean decision, a CDF, a marked record and a physical intervention are maintained throughout.

\section{Independent completed consequences}
\subsection{All factorially many assignments in a narrower physical model}
The independent file \texttt{next6\_decision\_theory/sorting\_intervention.tex} proves a global assignment rule for independent components with transferable reciprocal clock modules. Installing module $j$ at component $i$ must produce actual rates $r_jq_i$ and $r_j(1-q_i)$. Pairing larger $q_i$ with larger $r_j$ maximizes the equilibrium-safe mean to the all-failed state among all $N!$ assignments. The opposite order minimizes it. The exchange identity
\[
 [q_i+(1-q_i)e^{-at}][q_j+(1-q_j)e^{-bt}]
 -[q_i+(1-q_i)e^{-bt}][q_j+(1-q_j)e^{-at}]
 =(q_i-q_j)(e^{-bt}-e^{-at})
\]
proves the pointwise return comparison. This exact ordering survives an arbitrary independent ergodic spectator process with a specified singleton target; its preparation costs remain.

A dimension-independent bound on the logarithmic sensitivity of the nonrare return prefactor means constant, decision-level probability accuracy suffices. The completed acquisition uses $O(N\epsilon^{-2}\log(N/\delta))$ local reads at fixed margins, plus priced holds, timers and installation, followed by $O(N\log N)$ sorting. It does not need $1/N$ parameter accuracy or an absolute lifetime. Classical rearrangement and reliability-allocation precedents are credited; no general clock-module transport is inferred from biased calibration.

\subsection{Continuous design for desired completion}
The separate \texttt{next6\_direct\_decision} construction minimizes a desired rare completion time over a rational polytope of certified reciprocal kinetic multipliers. Its normalized objective $pR M(c)$ is convex with rarity-free bounds and curvature. A finite scalar oracle and inexact Frank--Wolfe optimize the continuum without a control grid or global matrix. This is completion acceleration, not maximizing a convex prevention objective. Bounded multipliers improve actual completion by at most their bounded factor.

Its evaluated elementary sufficient schedule at $N=128$, accuracy $0.1$, uses about $4.12\,10^8$ scalar calls and $10^{61.65}$ potential updates. These were not executed and are not lower bounds. Product-reference executions have a stronger ordinary solver and gains only about $1.10$--$1.23$. This useful global-design deduction was therefore not selected as the main capability.

\subsection{Quantum export without Hamiltonian identification}
The independent \texttt{next6\_physical/CONTROL\_FOR\_CALIBRATION.md} adds calibrated local source rotations to the accepted pure coherent Weyl ports. Covariant source/output frame correction and fresh Pauli averaging compare the full retained dilation, not only the reduced source. For unknown predetermined $H(t)$ with $\|H(t)\|\le\Lambda$, its finite export error obeys
\[
 2Ne^{-rT}+\frac{2\Lambda^2T^2}{B}
       +2(\Lambda+Nr)W_{\rm lost}+\epsilon_{\rm interface}.
\]
It works for every $r>0$ without Hamiltonian-learning exposures, with finite smooth pulses and charged lost output intervals. The random key is corrected and excluded from the returned target; retaining it does not have the same quadratic guarantee. Pure ports, large calibrated coherent receivers, protected memory, source strength and slew remain substantive. At $N=64$, $\Lambda=N$, $r=1$, $\epsilon=0.1$, a sufficient recipe uses about $23.9$ million blocks and peak per-site control norm $7.69\,10^{11}$ before receiver costs. Established randomized decoupling is credited. This is an additional-access resource trade, not same-access dominance over output-only acquisition or an economical receiver.

\section{Preservation and scientific status}
All 229 files of the preceding acquired-first-event package are preserved byte for byte, including its complete fitted-model and experiment-fed laws, operating and static coin constructions, confidence/preparation bills, independent results and nested earlier sources. The new archive's preservation receipt checks the original workspace, copied originals and ZIP members against the initial SHA-256 baseline. Historical failed or incomplete attempts remain historical files; the newest governing reports specify which interfaces were subsequently completed.

The result earned here is a complete conditional theorem, finite acquisition/decision/deployment construction, explicit same-access forecast separation, uniform safe-start mean comparison, and growing-family numerical execution. It is not evidence that a natural system satisfies the model, a prototype, a universal prevention optimizer, or demonstrated civilization-scale impact. The scientific principle is precise: when interventions share an invariant rare-event scale, the local kinetic information needed to decide can be much smaller than the information needed to forecast that scale. In the certified boundary regime even the initial full law can be discarded for this objective, because the theorem controls the mean uniformly over safe states.

\begin{thebibliography}{20}
\bibitem{aldousfill} D. Aldous and J. A. Fill, \emph{Reversible Markov Chains and Random Walks on Graphs}, unfinished monograph, 2002, recompiled 2014. Chapters 2--4. \url{https://www.stat.berkeley.edu/~aldous/RWG/book.html}.
\bibitem{bovier} A. Bovier, ``Metastability: a potential theoretic approach,'' \emph{Proceedings of the International Congress of Mathematicians}, 2006, vol. III, 499--518. DOI: 10.4171/022-3/26.
\bibitem{aldousbrown} D. Aldous and M. Brown, ``Inequalities for rare events in time-reversible Markov chains I,'' in \emph{Stochastic Inequalities}, IMS Lecture Notes--Monograph Series 22, 1992, 1--16; part II, \emph{Stochastic Processes and their Applications} 44, 1993, 15--25.
\bibitem{loper} J. Loper, G. Zhou and S. Geman, ``Capacities and efficient computation of first-passage probabilities,'' \emph{Physical Review E} 102, 2020, 023304. Full primary paper: \url{https://stanniszhou.github.io/papers/diffusion_manuscript.pdf}.
\bibitem{lecam} E. Kaufmann, O. Cappe and A. Garivier, ``On the complexity of best-arm identification in multi-armed bandit models,'' \emph{Journal of Machine Learning Research} 17, 2016.
\bibitem{hayes} T. P. Hayes, ``A simple condition implying rapid mixing of single-site dynamics on spin systems,'' \emph{FOCS}, 2006, 39--46.
\bibitem{nacu} S. Nacu and Y. Peres, ``Fast simulation of new coins from old,'' \emph{Annals of Applied Probability} 15, 2005, 93--115.
\bibitem{astra} Accepted Astra constructions, exact unchanged source tree \texttt{preserved\_acquired\_first\_event/} in the accompanying archive. Governing public and local source registers, proofs, programs, failures and limits are preserved byte for byte.
\end{thebibliography}
\end{document}
